%%%%%%%%%%%%%%%%%%%%%%%%%%%%%%%%%%%%%%%%%%%%%%%%%%%%%%%%%%%%%%%%%%%%%%%%%%%%%%%%%%%%%
% BỘ 16 CÂU HỎI — "PHẦN KHÔNG BỊ GẠCH LÀ MIỀN NGHIỆM CỦA BẤT PHƯƠNG TRÌNH NÀO?"
% Chủ đề: Bất phương trình bậc nhất hai ẩn (Toán 10 — Chương II, Bài 1)
% Chuẩn: ex_test v3.3.4 + preamble cautrucDT (xem CheatSheet-LaTeX-VNPmath.md)
%
% SẮP XẾP THEO 4 DẠNG ĐƯỜNG THẲNG BỜ:
%   Dạng 1. Bờ là đường thẳng x = m            (5 câu: 1 → 5)   [m = -2, -1, 0, 1, 2]
%   Dạng 2. Bờ là đường thẳng y = m            (4 câu: 6 → 9)   [m = -1, 0, 1, 2]
%   Dạng 3. Bờ là đường thẳng ax+by = c, ab<0  (4 câu: 10 → 13) — hệ số góc dương
%   Dạng 4. Bờ là đường thẳng ax+by = c, ab>0  (3 câu: 14 → 16) — hệ số góc âm
%
% QUY ƯỚC VẼ HÌNH (áp dụng cho cả 16 câu):
%   1) Khung hình [xmin;xmax] × [ymin;ymax] khai báo trong declare function,
%      cùng hàm đường bờ (với bờ xiên) hoặc tham số m (với bờ x=m, y=m).
%   2) Tô theo đúng ví dụ:
%        \fill[pattern=north east lines] (xmin,ymin) rectangle (xmax,ymax);   % tô kín khung
%        \fill[white] <đa giác MIỀN NGHIỆM>;                                  % xoá nét ở miền nghiệm
%      ⇒ phần nhìn thấy bị gạch = nửa mặt phẳng KHÔNG phải miền nghiệm (như các hình trước).
%   3) CHỌN KHUNG SAO CHO ĐƯỜNG BỜ d ĐI QUA TÂM KHUNG ⇒ phần gạch chiếm ĐÚNG 50% diện tích.
%      Công thức: tâm khung = ((xmin+xmax)/2 ; (ymin+ymax)/2) phải là một điểm thuộc d.
%   4) Bờ dạng \geq / \leq  → nét liền:  \draw[thick] ...
%      Bờ dạng > / <        → nét đứt:  \draw[thick,dashed] ...
%      Bờ trùng trục toạ độ (câu 1 và câu 6) thì không vẽ thêm đường bờ.
%   5) Bờ xiên: bọc vùng tô và nét bờ trong \begin{scope}\clip (xmin,ymin) rectangle (xmax,ymax);
%      ...\end{scope} để nét bờ và mảng tô không tràn ra ngoài khung.
%   6) Lưới trục: vẽ trục SAU khi tô (để số trên trục không bị nét gạch cắt qua);
%      chỉ ghi số tại các giá trị nguyên nằm trong khung (bỏ 0, ký hiệu bằng $O$).
%
% BẢNG ĐÁP ÁN (đối chiếu nhanh):
%   1.B  2.A  3.C  4.D  5.D  6.D  7.A  8.C
%   9.B  10.D 11.D 12.A 13.B 14.C 15.B 16.C
%
% CÁCH DÙNG
% 1) File này tự mở/đóng file đáp án riêng: ans/ans-0D2-B1-MienNghiem.
%    → Trong input10.tex chỉ cần thêm 1 dòng:
%        %%%%%%%%%%%%%%%%%%%%%%%%%%%%%%%%%%%%%%%%%%%%%%%%%%%%%%%%%%%%%%%%%%%%%%%%%%%%%%%%%%%%%
% BỘ 16 CÂU HỎI — "PHẦN KHÔNG BỊ GẠCH LÀ MIỀN NGHIỆM CỦA BẤT PHƯƠNG TRÌNH NÀO?"
% Chủ đề: Bất phương trình bậc nhất hai ẩn (Toán 10 — Chương II, Bài 1)
% Chuẩn: ex_test v3.3.4 + preamble cautrucDT (xem CheatSheet-LaTeX-VNPmath.md)
%
% SẮP XẾP THEO 4 DẠNG ĐƯỜNG THẲNG BỜ:
%   Dạng 1. Bờ là đường thẳng x = m            (5 câu: 1 → 5)   [m = -2, -1, 0, 1, 2]
%   Dạng 2. Bờ là đường thẳng y = m            (4 câu: 6 → 9)   [m = -1, 0, 1, 2]
%   Dạng 3. Bờ là đường thẳng ax+by = c, ab<0  (4 câu: 10 → 13) — hệ số góc dương
%   Dạng 4. Bờ là đường thẳng ax+by = c, ab>0  (3 câu: 14 → 16) — hệ số góc âm
%
% QUY ƯỚC VẼ HÌNH (áp dụng cho cả 16 câu):
%   1) Khung hình [xmin;xmax] × [ymin;ymax] khai báo trong declare function,
%      cùng hàm đường bờ (với bờ xiên) hoặc tham số m (với bờ x=m, y=m).
%   2) Tô theo đúng ví dụ:
%        \fill[pattern=north east lines] (xmin,ymin) rectangle (xmax,ymax);   % tô kín khung
%        \fill[white] <đa giác MIỀN NGHIỆM>;                                  % xoá nét ở miền nghiệm
%      ⇒ phần nhìn thấy bị gạch = nửa mặt phẳng KHÔNG phải miền nghiệm (như các hình trước).
%   3) CHỌN KHUNG SAO CHO ĐƯỜNG BỜ d ĐI QUA TÂM KHUNG ⇒ phần gạch chiếm ĐÚNG 50% diện tích.
%      Công thức: tâm khung = ((xmin+xmax)/2 ; (ymin+ymax)/2) phải là một điểm thuộc d.
%   4) Bờ dạng \geq / \leq  → nét liền:  \draw[thick] ...
%      Bờ dạng > / <        → nét đứt:  \draw[thick,dashed] ...
%      Bờ trùng trục toạ độ (câu 1 và câu 6) thì không vẽ thêm đường bờ.
%   5) Bờ xiên: bọc vùng tô và nét bờ trong \begin{scope}\clip (xmin,ymin) rectangle (xmax,ymax);
%      ...\end{scope} để nét bờ và mảng tô không tràn ra ngoài khung.
%   6) Lưới trục: vẽ trục SAU khi tô (để số trên trục không bị nét gạch cắt qua);
%      chỉ ghi số tại các giá trị nguyên nằm trong khung (bỏ 0, ký hiệu bằng $O$).
%
% BẢNG ĐÁP ÁN (đối chiếu nhanh):
%   1.B  2.A  3.C  4.D  5.D  6.D  7.A  8.C
%   9.B  10.D 11.D 12.A 13.B 14.C 15.B 16.C
%
% CÁCH DÙNG
% 1) File này tự mở/đóng file đáp án riêng: ans/ans-0D2-B1-MienNghiem.
%    → Trong input10.tex chỉ cần thêm 1 dòng:
%        %%%%%%%%%%%%%%%%%%%%%%%%%%%%%%%%%%%%%%%%%%%%%%%%%%%%%%%%%%%%%%%%%%%%%%%%%%%%%%%%%%%%%
% BỘ 16 CÂU HỎI — "PHẦN KHÔNG BỊ GẠCH LÀ MIỀN NGHIỆM CỦA BẤT PHƯƠNG TRÌNH NÀO?"
% Chủ đề: Bất phương trình bậc nhất hai ẩn (Toán 10 — Chương II, Bài 1)
% Chuẩn: ex_test v3.3.4 + preamble cautrucDT (xem CheatSheet-LaTeX-VNPmath.md)
%
% SẮP XẾP THEO 4 DẠNG ĐƯỜNG THẲNG BỜ:
%   Dạng 1. Bờ là đường thẳng x = m            (5 câu: 1 → 5)   [m = -2, -1, 0, 1, 2]
%   Dạng 2. Bờ là đường thẳng y = m            (4 câu: 6 → 9)   [m = -1, 0, 1, 2]
%   Dạng 3. Bờ là đường thẳng ax+by = c, ab<0  (4 câu: 10 → 13) — hệ số góc dương
%   Dạng 4. Bờ là đường thẳng ax+by = c, ab>0  (3 câu: 14 → 16) — hệ số góc âm
%
% QUY ƯỚC VẼ HÌNH (áp dụng cho cả 16 câu):
%   1) Khung hình [xmin;xmax] × [ymin;ymax] khai báo trong declare function,
%      cùng hàm đường bờ (với bờ xiên) hoặc tham số m (với bờ x=m, y=m).
%   2) Tô theo đúng ví dụ:
%        \fill[pattern=north east lines] (xmin,ymin) rectangle (xmax,ymax);   % tô kín khung
%        \fill[white] <đa giác MIỀN NGHIỆM>;                                  % xoá nét ở miền nghiệm
%      ⇒ phần nhìn thấy bị gạch = nửa mặt phẳng KHÔNG phải miền nghiệm (như các hình trước).
%   3) CHỌN KHUNG SAO CHO ĐƯỜNG BỜ d ĐI QUA TÂM KHUNG ⇒ phần gạch chiếm ĐÚNG 50% diện tích.
%      Công thức: tâm khung = ((xmin+xmax)/2 ; (ymin+ymax)/2) phải là một điểm thuộc d.
%   4) Bờ dạng \geq / \leq  → nét liền:  \draw[thick] ...
%      Bờ dạng > / <        → nét đứt:  \draw[thick,dashed] ...
%      Bờ trùng trục toạ độ (câu 1 và câu 6) thì không vẽ thêm đường bờ.
%   5) Bờ xiên: bọc vùng tô và nét bờ trong \begin{scope}\clip (xmin,ymin) rectangle (xmax,ymax);
%      ...\end{scope} để nét bờ và mảng tô không tràn ra ngoài khung.
%   6) Lưới trục: vẽ trục SAU khi tô (để số trên trục không bị nét gạch cắt qua);
%      chỉ ghi số tại các giá trị nguyên nằm trong khung (bỏ 0, ký hiệu bằng $O$).
%
% BẢNG ĐÁP ÁN (đối chiếu nhanh):
%   1.B  2.A  3.C  4.D  5.D  6.D  7.A  8.C
%   9.B  10.D 11.D 12.A 13.B 14.C 15.B 16.C
%
% CÁCH DÙNG
% 1) File này tự mở/đóng file đáp án riêng: ans/ans-0D2-B1-MienNghiem.
%    → Trong input10.tex chỉ cần thêm 1 dòng:
%        %%%%%%%%%%%%%%%%%%%%%%%%%%%%%%%%%%%%%%%%%%%%%%%%%%%%%%%%%%%%%%%%%%%%%%%%%%%%%%%%%%%%%
% BỘ 16 CÂU HỎI — "PHẦN KHÔNG BỊ GẠCH LÀ MIỀN NGHIỆM CỦA BẤT PHƯƠNG TRÌNH NÀO?"
% Chủ đề: Bất phương trình bậc nhất hai ẩn (Toán 10 — Chương II, Bài 1)
% Chuẩn: ex_test v3.3.4 + preamble cautrucDT (xem CheatSheet-LaTeX-VNPmath.md)
%
% SẮP XẾP THEO 4 DẠNG ĐƯỜNG THẲNG BỜ:
%   Dạng 1. Bờ là đường thẳng x = m            (5 câu: 1 → 5)   [m = -2, -1, 0, 1, 2]
%   Dạng 2. Bờ là đường thẳng y = m            (4 câu: 6 → 9)   [m = -1, 0, 1, 2]
%   Dạng 3. Bờ là đường thẳng ax+by = c, ab<0  (4 câu: 10 → 13) — hệ số góc dương
%   Dạng 4. Bờ là đường thẳng ax+by = c, ab>0  (3 câu: 14 → 16) — hệ số góc âm
%
% QUY ƯỚC VẼ HÌNH (áp dụng cho cả 16 câu):
%   1) Khung hình [xmin;xmax] × [ymin;ymax] khai báo trong declare function,
%      cùng hàm đường bờ (với bờ xiên) hoặc tham số m (với bờ x=m, y=m).
%   2) Tô theo đúng ví dụ:
%        \fill[pattern=north east lines] (xmin,ymin) rectangle (xmax,ymax);   % tô kín khung
%        \fill[white] <đa giác MIỀN NGHIỆM>;                                  % xoá nét ở miền nghiệm
%      ⇒ phần nhìn thấy bị gạch = nửa mặt phẳng KHÔNG phải miền nghiệm (như các hình trước).
%   3) CHỌN KHUNG SAO CHO ĐƯỜNG BỜ d ĐI QUA TÂM KHUNG ⇒ phần gạch chiếm ĐÚNG 50% diện tích.
%      Công thức: tâm khung = ((xmin+xmax)/2 ; (ymin+ymax)/2) phải là một điểm thuộc d.
%   4) Bờ dạng \geq / \leq  → nét liền:  \draw[thick] ...
%      Bờ dạng > / <        → nét đứt:  \draw[thick,dashed] ...
%      Bờ trùng trục toạ độ (câu 1 và câu 6) thì không vẽ thêm đường bờ.
%   5) Bờ xiên: bọc vùng tô và nét bờ trong \begin{scope}\clip (xmin,ymin) rectangle (xmax,ymax);
%      ...\end{scope} để nét bờ và mảng tô không tràn ra ngoài khung.
%   6) Lưới trục: vẽ trục SAU khi tô (để số trên trục không bị nét gạch cắt qua);
%      chỉ ghi số tại các giá trị nguyên nằm trong khung (bỏ 0, ký hiệu bằng $O$).
%
% BẢNG ĐÁP ÁN (đối chiếu nhanh):
%   1.B  2.A  3.C  4.D  5.D  6.D  7.A  8.C
%   9.B  10.D 11.D 12.A 13.B 14.C 15.B 16.C
%
% CÁCH DÙNG
% 1) File này tự mở/đóng file đáp án riêng: ans/ans-0D2-B1-MienNghiem.
%    → Trong input10.tex chỉ cần thêm 1 dòng:
%        \input{data/10/0D2/TN-0D2-B1-MienNghiem.tex}
% 2) Nếu muốn GỘP vào TN-0D2-B1.tex: xoá 2 dòng \Opensolutionfile/\Closesolutionfile
%    ở file này rồi dán phần các câu hỏi vào TRƯỚC dòng \Closesolutionfile{ans} của file đó.
% 3) In bảng đáp án: \inputansbox{10}{ans/ans-0D2-B1-MienNghiem}
% 4) Lời giải ngắn đã điền sẵn trong \loigiai{}; xoá nội dung nếu muốn để trống như mẫu gốc.
%    QUY ƯỚC LỜI GIẢI: chỉ phân tích miền KHÔNG bị gạch — nêu phương trình đường bờ (kèm ghi
%    chú nét liền/nét đứt), rồi lấy điểm $M$ thuộc phần không bị gạch (ưu tiên điểm có ít nhất
%    một toạ độ bằng 0) thay vào bất phương trình để kết luận.
%    - Nếu lấy gốc toạ độ thì gọi là $O(0;0)$ (KHÔNG dùng $M(0;0)$).
% 5) Đang dùng \subsubsection làm tiêu đề dạng (đánh số 1., 2., 3., 4.; KHÔNG reset số câu,
%    nên bảng đáp án vẫn liên tục 1 → 16). Nếu muốn khung "DẠNG TOÁN" thì thay bằng
%    \begin{dang}{...} ... \end{dang} — lưu ý môi trường dang tự reset số câu về 1.
% 6) Quy ước nhãn câu: %[Nguồn]%[Người biên soạn]%[Mã ID].
%    Mã ID dùng ở đây: 0D2N1-2 (nhận biết — bờ song song trục) và 0D2H1-2 (thông hiểu — bờ xiên),
%    khớp với data/10/0D2/0-KNTT-C2B3-X.tex.
%
% GHI CHÚ: hai câu bờ x = m, y = m dạng nét đứt trước đây dùng m = 3 (khung phải lệch, khó đạt
% 50% mà vẫn thấy gốc O) đã đổi thành m = -2 (câu 5) và m = -1 (câu 9) để khung đối xứng qua d.
%%%%%%%%%%%%%%%%%%%%%%%%%%%%%%%%%%%%%%%%%%%%%%%%%%%%%%%%%%%%%%%%%%%%%%%%%%%%%%%%%%%%%

\Opensolutionfile{ans}[ans/ans-0D2-B1-MienNghiem]

%%%%%%===================================================================%%%%%%
%%%%%%============================= DẠNG 1 ===============================%%%%%%
%%%%%% Bờ là đường thẳng song song với trục tung: x = m
%%%%%%===================================================================%%%%%%
\subsubsection{Bờ là đường thẳng $x=m$}

%%%==============Câu_1 — x = 0 (trục tung), tô bên trái, bờ NÉT LIỀN=============%%%
\begin{ex}%[VNPmath - tự biên soạn]%[GV. Vũ Ngọc Phát]%[0D2N1-2]
	\immini[thm]{Phần không bị gạch trong hình vẽ là miền nghiệm của bất phương trình nào sau đây?
		\choice
		{$x\leq 0$}
		{\True $x\geq 0$}
		{$y\leq 0$}
		{$y\geq 0$}}
	{\begin{tikzpicture}[scale=.7,thick,>=stealth,declare function={xmin=-3;xmax=3;ymin=-2;ymax=3;m=0;}]
			\fill[pattern=north east lines] (xmin,ymin) rectangle (xmax,ymax);
			\fill[white] (m,ymin) rectangle (xmax,ymax);
			\draw[->] (xmin,0) -- (xmax,0)node[above]{$x$};
			\foreach \x in {-2,-1,1,2,3}
			\draw[shift={(\x,0)}] (0pt,2pt) -- (0pt,-2pt) node[below] {\footnotesize $\x$};
			\draw[->,color=black] (0,ymin) -- (0,ymax)node[left]{$y$};
			\foreach \y in {-1,1,2}
			\draw[shift={(0,\y)}] (2pt,0pt) -- (-2pt,0pt) node[left] {\footnotesize $\y$};
			\node[above left] at (0,0){$O$};
		\end{tikzpicture}}
	\loigiai{
		Đường thẳng bờ là trục tung $d\colon x=0$ (vẽ nét liền nên kể cả bờ).\\
		Lấy điểm $M(1;0)$ thuộc phần không bị gạch, ta có $1\geq 0$ nên phần không bị gạch là miền nghiệm của bất phương trình $x\geq 0$.
	}
\end{ex}

%%%==============Câu_2 — x = 1, tô bên trái, bờ NÉT LIỀN=============%%%
\begin{ex}%[VNPmath - tự biên soạn]%[GV. Vũ Ngọc Phát]%[0D2N1-2]
	\immini[thm]{Phần không bị gạch trong hình vẽ là miền nghiệm của bất phương trình nào sau đây?
		\choice
		{\True $x\geq 1$}
		{$x\leq 1$}
		{$x\geq -1$}
		{$x\leq -1$}}
	{\begin{tikzpicture}[scale=.7,thick,>=stealth,declare function={xmin=-2;xmax=4;ymin=-2;ymax=3;m=1;}]
			\fill[pattern=north east lines] (xmin,ymin) rectangle (xmax,ymax);
			\fill[white] (m,ymin) rectangle (xmax,ymax);
			\draw[thick] (m,ymin) -- (m,ymax);
			\draw[->] (xmin,0) -- (xmax,0)node[above]{$x$};
			\foreach \x in {-1,1,2,3,4}
			\draw[shift={(\x,0)}] (0pt,2pt) -- (0pt,-2pt) node[below] {\footnotesize $\x$};
			\draw[->,color=black] (0,ymin) -- (0,ymax)node[left]{$y$};
			\foreach \y in {-1,1,2}
			\draw[shift={(0,\y)}] (2pt,0pt) -- (-2pt,0pt) node[left] {\footnotesize $\y$};
			\node[above left] at (0,0){$O$};
		\end{tikzpicture}}
	\loigiai{
		Đường thẳng bờ $d\colon x=1$ (vẽ nét liền nên kể cả bờ).\\
		Lấy điểm $M(2;0)$ thuộc phần không bị gạch, ta có $2\geq 1$ nên phần không bị gạch là miền nghiệm của bất phương trình $x\geq 1$.
	}
\end{ex}

%%%==============Câu_3 — x = -1, tô bên phải, bờ NÉT LIỀN=============%%%
\begin{ex}%[VNPmath - tự biên soạn]%[GV. Vũ Ngọc Phát]%[0D2N1-2]
	\immini[thm]{Phần không bị gạch trong hình vẽ là miền nghiệm của bất phương trình nào sau đây?
		\choice
		{$x\geq -1$}
		{$x\geq 1$}
		{\True $x\leq -1$}
		{$x\leq 1$}}
	{\begin{tikzpicture}[scale=.7,thick,>=stealth,declare function={xmin=-4;xmax=2;ymin=-2;ymax=3;m=-1;}]
			\fill[pattern=north east lines] (xmin,ymin) rectangle (xmax,ymax);
			\fill[white] (xmin,ymin) rectangle (m,ymax);
			\draw[thick] (m,ymin) -- (m,ymax);
			\draw[->] (xmin,0) -- (xmax,0)node[above]{$x$};
			\foreach \x in {-3,-2,-1,1,2}
			\draw[shift={(\x,0)}] (0pt,2pt) -- (0pt,-2pt) node[below] {\footnotesize $\x$};
			\draw[->,color=black] (0,ymin) -- (0,ymax)node[left]{$y$};
			\foreach \y in {-1,1,2}
			\draw[shift={(0,\y)}] (2pt,0pt) -- (-2pt,0pt) node[left] {\footnotesize $\y$};
			\node[above left] at (0,0){$O$};
		\end{tikzpicture}}
	\loigiai{
		Đường thẳng bờ $d\colon x=-1$ (vẽ nét liền nên kể cả bờ).\\
		Lấy điểm $M(-2;0)$ thuộc phần không bị gạch, ta có $-2\leq -1$ nên phần không bị gạch là miền nghiệm của bất phương trình $x\leq -1$.
	}
\end{ex}

%%%==============Câu_4 — x = 2, tô bên phải, bờ NÉT ĐỨT=============%%%
\begin{ex}%[VNPmath - tự biên soạn]%[GV. Vũ Ngọc Phát]%[0D2N1-2]
	\immini[thm]{Phần không bị gạch (không kể đường thẳng $d$) trong hình vẽ là miền nghiệm của bất phương trình nào sau đây?
		\choice
		{$x>2$}
		{$x<3$}
		{$x>3$}
		{\True $x<2$}}
	{\begin{tikzpicture}[scale=.7,thick,>=stealth,declare function={xmin=-1;xmax=5;ymin=-2;ymax=3;m=2;}]
			\fill[pattern=north east lines] (xmin,ymin) rectangle (xmax,ymax);
			\fill[white] (xmin,ymin) rectangle (m,ymax);
			\draw[thick,dashed] (m,ymin) -- (m,ymax);
			\draw[->] (xmin,0) -- (xmax,0)node[above]{$x$};
			\foreach \x in {-1,1,2,3,4,5}
			\draw[shift={(\x,0)}] (0pt,2pt) -- (0pt,-2pt) node[below] {\footnotesize $\x$};
			\draw[->,color=black] (0,ymin) -- (0,ymax)node[left]{$y$};
			\foreach \y in {-1,1,2}
			\draw[shift={(0,\y)}] (2pt,0pt) -- (-2pt,0pt) node[left] {\footnotesize $\y$};
			\node[above left] at (0,0){$O$};
		\end{tikzpicture}}
	\loigiai{
		Đường thẳng bờ $d\colon x=2$ vẽ nét đứt nên không kể bờ.\\
		Lấy điểm $O(0;0)$ thuộc phần không bị gạch, ta có $0<2$ nên phần không bị gạch là miền nghiệm của bất phương trình $x<2$.
	}
\end{ex}

%%%==============Câu_5 — x = -2, tô bên trái, bờ NÉT ĐỨT=============%%%
\begin{ex}%[VNPmath - tự biên soạn]%[GV. Vũ Ngọc Phát]%[0D2N1-2]
	\immini[thm]{Phần không bị gạch (không kể đường thẳng $d$) trong hình vẽ là miền nghiệm của bất phương trình nào sau đây?
		\choice
		{$x<-3$}
		{$x>-3$}
		{$x<-2$}
		{\True $x>-2$}}
	{\begin{tikzpicture}[scale=.7,thick,>=stealth,declare function={xmin=-5;xmax=1;ymin=-2;ymax=3;m=-2;}]
			\fill[pattern=north east lines] (xmin,ymin) rectangle (xmax,ymax);
			\fill[white] (m,ymin) rectangle (xmax,ymax);
			\draw[thick,dashed] (m,ymin) -- (m,ymax);
			\draw[->] (xmin,0) -- (xmax,0)node[above]{$x$};
			\foreach \x in {-4,-3,-2,-1,1}
			\draw[shift={(\x,0)}] (0pt,2pt) -- (0pt,-2pt) node[below] {\footnotesize $\x$};
			\draw[->,color=black] (0,ymin) -- (0,ymax)node[left]{$y$};
			\foreach \y in {-1,1,2}
			\draw[shift={(0,\y)}] (2pt,0pt) -- (-2pt,0pt) node[left] {\footnotesize $\y$};
			\node[above left] at (0,0){$O$};
		\end{tikzpicture}}
	\loigiai{
		Đường thẳng bờ $d\colon x=-2$ vẽ nét đứt nên không kể bờ.\\
		Lấy điểm $O(0;0)$ thuộc phần không bị gạch, ta có $0>-2$ nên phần không bị gạch là miền nghiệm của bất phương trình $x>-2$.
	}
\end{ex}

%%%%%%===================================================================%%%%%%
%%%%%%============================= DẠNG 2 ===============================%%%%%%
%%%%%% Bờ là đường thẳng song song với trục hoành: y = m
%%%%%%===================================================================%%%%%%
\subsubsection{Bờ là đường thẳng $y=m$}

%%%==============Câu_6 — y = 0 (trục hoành), tô phía dưới, bờ NÉT LIỀN=============%%%
\begin{ex}%[VNPmath - tự biên soạn]%[GV. Vũ Ngọc Phát]%[0D2N1-2]
	\immini[thm]{Phần không bị gạch trong hình vẽ là miền nghiệm của bất phương trình nào sau đây?
		\choice
		{$y\leq 0$}
		{$x\leq 0$}
		{$x\geq 0$}
		{\True $y\geq 0$}}
	{\begin{tikzpicture}[scale=.7,thick,>=stealth,declare function={xmin=-2;xmax=4;ymin=-2.5;ymax=2.5;m=0;}]
			\fill[pattern=north east lines] (xmin,ymin) rectangle (xmax,ymax);
			\fill[white] (xmin,m) rectangle (xmax,ymax);
			\draw[->] (xmin,0) -- (xmax,0)node[above]{$x$};
			\foreach \x in {-1,1,2,3,4}
			\draw[shift={(\x,0)}] (0pt,2pt) -- (0pt,-2pt) node[below] {\footnotesize $\x$};
			\draw[->,color=black] (0,ymin) -- (0,ymax)node[left]{$y$};
			\foreach \y in {-2,-1,1,2}
			\draw[shift={(0,\y)}] (2pt,0pt) -- (-2pt,0pt) node[left] {\footnotesize $\y$};
			\node[above left] at (0,0){$O$};
		\end{tikzpicture}}
	\loigiai{
		Đường thẳng bờ là trục hoành $d\colon y=0$ (vẽ nét liền nên kể cả bờ).\\
		Lấy điểm $M(0;1)$ thuộc phần không bị gạch, ta có $1\geq 0$ nên phần không bị gạch là miền nghiệm của bất phương trình $y\geq 0$.
	}
\end{ex}

%%%==============Câu_7 — y = 1, tô phía dưới, bờ NÉT ĐỨT=============%%%
\begin{ex}%[VNPmath - tự biên soạn]%[GV. Vũ Ngọc Phát]%[0D2N1-2]
	\immini[thm]{Phần không bị gạch (không kể đường thẳng $d$) trong hình vẽ là miền nghiệm của bất phương trình nào sau đây?
		\choice
		{\True $y>1$}
		{$y<1$}
		{$y>2$}
		{$y<2$}}
	{\begin{tikzpicture}[scale=.7,thick,>=stealth,declare function={xmin=-2;xmax=4;ymin=-1.5;ymax=3.5;m=1;}]
			\fill[pattern=north east lines] (xmin,ymin) rectangle (xmax,ymax);
			\fill[white] (xmin,m) rectangle (xmax,ymax);
			\draw[thick,dashed] (xmin,m) -- (xmax,m);
			\draw[->] (xmin,0) -- (xmax,0)node[above]{$x$};
			\foreach \x in {-1,1,2,3,4}
			\draw[shift={(\x,0)}] (0pt,2pt) -- (0pt,-2pt) node[below] {\footnotesize $\x$};
			\draw[->,color=black] (0,ymin) -- (0,ymax)node[left]{$y$};
			\foreach \y in {-1,1,2,3}
			\draw[shift={(0,\y)}] (2pt,0pt) -- (-2pt,0pt) node[left] {\footnotesize $\y$};
			\node[above left] at (0,0){$O$};
		\end{tikzpicture}}
	\loigiai{
		Đường thẳng bờ $d\colon y=1$ vẽ nét đứt nên không kể bờ.\\
		Lấy điểm $M(0;2)$ thuộc phần không bị gạch, ta có $2>1$ nên phần không bị gạch là miền nghiệm của bất phương trình $y>1$.
	}
\end{ex}

%%%==============Câu_8 — y = 2, tô phía trên, bờ NÉT LIỀN=============%%%
\begin{ex}%[VNPmath - tự biên soạn]%[GV. Vũ Ngọc Phát]%[0D2N1-2]
	\immini[thm]{Phần không bị gạch trong hình vẽ là miền nghiệm của bất phương trình nào sau đây?
		\choice
		{$y\geq 2$}
		{$y\geq -2$}
		{\True $y\leq 2$}
		{$y\leq -2$}}
	{\begin{tikzpicture}[scale=.7,thick,>=stealth,declare function={xmin=-2;xmax=4;ymin=-0.5;ymax=4.5;m=2;}]
			\fill[pattern=north east lines] (xmin,ymin) rectangle (xmax,ymax);
			\fill[white] (xmin,ymin) rectangle (xmax,m);
			\draw[thick] (xmin,m) -- (xmax,m);
			\draw[->] (xmin,0) -- (xmax,0)node[above]{$x$};
			\foreach \x in {-1,1,2,3,4}
			\draw[shift={(\x,0)}] (0pt,2pt) -- (0pt,-2pt) node[below] {\footnotesize $\x$};
			\draw[->,color=black] (0,ymin) -- (0,ymax)node[left]{$y$};
			\foreach \y in {1,2,3,4}
			\draw[shift={(0,\y)}] (2pt,0pt) -- (-2pt,0pt) node[left] {\footnotesize $\y$};
			\node[above left] at (0,0){$O$};
		\end{tikzpicture}}
	\loigiai{
		Đường thẳng bờ $d\colon y=2$ (vẽ nét liền nên kể cả bờ).\\
		Lấy điểm $O(0;0)$ thuộc phần không bị gạch, ta có $0\leq 2$ nên phần không bị gạch là miền nghiệm của bất phương trình $y\leq 2$.
	}
\end{ex}

%%%==============Câu_9 — y = -1, tô phía trên, bờ NÉT ĐỨT=============%%%
\begin{ex}%[VNPmath - tự biên soạn]%[GV. Vũ Ngọc Phát]%[0D2N1-2]
	\immini[thm]{Phần không bị gạch (không kể đường thẳng $d$) trong hình vẽ là miền nghiệm của bất phương trình nào sau đây?
		\choice
		{$y>-1$}
		{\True $y<-1$}
		{$y>-2$}
		{$y<-2$}}
	{\begin{tikzpicture}[scale=.7,thick,>=stealth,declare function={xmin=-2;xmax=4;ymin=-3.5;ymax=1.5;m=-1;}]
			\fill[pattern=north east lines] (xmin,ymin) rectangle (xmax,ymax);
			\fill[white] (xmin,ymin) rectangle (xmax,m);
			\draw[thick,dashed] (xmin,m) -- (xmax,m);
			\draw[->] (xmin,0) -- (xmax,0)node[above]{$x$};
			\foreach \x in {-1,1,2,3,4}
			\draw[shift={(\x,0)}] (0pt,2pt) -- (0pt,-2pt) node[below] {\footnotesize $\x$};
			\draw[->,color=black] (0,ymin) -- (0,ymax)node[left]{$y$};
			\foreach \y in {-3,-2,-1,1}
			\draw[shift={(0,\y)}] (2pt,0pt) -- (-2pt,0pt) node[left] {\footnotesize $\y$};
			\node[above left] at (0,0){$O$};
		\end{tikzpicture}}
	\loigiai{
		Đường thẳng bờ $d\colon y=-1$ vẽ nét đứt nên không kể bờ.\\
		Lấy điểm $M(0;-2)$ thuộc phần không bị gạch, ta có $-2<-1$ nên phần không bị gạch là miền nghiệm của bất phương trình $y<-1$.
	}
\end{ex}

%%%%%%===================================================================%%%%%%
%%%%%%============================= DẠNG 3 ===============================%%%%%%
%%%%%% Bờ là đường thẳng ax+by = c với ab < 0  (hệ số góc dương)
%%%%%%===================================================================%%%%%%
\subsubsection{Bờ là đường thẳng $ax+by=c$ với $ab<0$}

%%%==============Câu_10 — y = x (x - y = 0), tô nửa trên - trái, bờ NÉT LIỀN=============%%%
\begin{ex}%[VNPmath - tự biên soạn]%[GV. Vũ Ngọc Phát]%[0D2H1-2]
	\immini[thm]{Phần không bị gạch trong hình vẽ là miền nghiệm của bất phương trình nào sau đây?
		\choice
		{$y\geq x$}
		{$y\geq -x$}
		{$y\leq -x$}
		{\True $y\leq x$}}
	{\begin{tikzpicture}[scale=.7,thick,>=stealth,declare function={xmin=-2.5;xmax=3.5;ymin=-2;ymax=3;f(\x)=\x;}]
			\begin{scope}
				\clip (xmin,ymin) rectangle (xmax,ymax);
				\fill[pattern=north east lines] (xmin,ymin) rectangle (xmax,ymax);
				\fill[white] (xmin,{f(xmin)}) -- (xmax,{f(xmax)}) -- (xmax,ymin) -- (xmin,ymin) -- cycle;
				\draw[thick] (xmin,{f(xmin)}) -- (xmax,{f(xmax)});
			\end{scope}
			\draw[->] (xmin,0) -- (xmax,0)node[above]{$x$};
			\foreach \x in {-2,-1,1,2,3}
			\draw[shift={(\x,0)}] (0pt,2pt) -- (0pt,-2pt) node[below] {\footnotesize $\x$};
			\draw[->,color=black] (0,ymin) -- (0,ymax)node[left]{$y$};
			\foreach \y in {-1,1,2}
			\draw[shift={(0,\y)}] (2pt,0pt) -- (-2pt,0pt) node[left] {\footnotesize $\y$};
			\node[above left] at (0,0){$O$};
		\end{tikzpicture}}
	\loigiai{
		Đường thẳng bờ đi qua gốc toạ độ $O$ và điểm $(1;1)$ nên có phương trình $y=x$ (vẽ nét liền nên kể cả bờ).\\
		Lấy điểm $M(1;0)$ thuộc phần không bị gạch, ta có $0\leq 1$ nên phần không bị gạch là miền nghiệm của bất phương trình $y\leq x$.
	}
\end{ex}

%%%==============Câu_11 — y = 2x (2x - y = 0), tô nửa trên - trái, bờ NÉT LIỀN=============%%%
\begin{ex}%[VNPmath - tự biên soạn]%[GV. Vũ Ngọc Phát]%[0D2H1-2]
	\immini[thm]{Phần không bị gạch trong hình vẽ là miền nghiệm của bất phương trình nào sau đây?
		\choice
		{$y\geq 2x$}
		{$2y\geq x$}
		{$2y\leq x$}
		{\True $y\leq 2x$}}
	{\begin{tikzpicture}[scale=.7,thick,>=stealth,declare function={xmin=-2.5;xmax=3.5;ymin=-1.5;ymax=3.5;f(\x)=2*(\x);}]
			\begin{scope}
				\clip (xmin,ymin) rectangle (xmax,ymax);
				\fill[pattern=north east lines] (xmin,ymin) rectangle (xmax,ymax);
				\fill[white] (xmin,{f(xmin)}) -- (xmax,{f(xmax)}) -- (xmax,ymin) -- (xmin,ymin) -- cycle;
				\draw[thick] (xmin,{f(xmin)}) -- (xmax,{f(xmax)});
			\end{scope}
			\draw[->] (xmin,0) -- (xmax,0)node[above]{$x$};
			\foreach \x in {-2,-1,1,2,3}
			\draw[shift={(\x,0)}] (0pt,2pt) -- (0pt,-2pt) node[below] {\footnotesize $\x$};
			\draw[->,color=black] (0,ymin) -- (0,ymax)node[left]{$y$};
			\foreach \y in {-1,1,2,3}
			\draw[shift={(0,\y)}] (2pt,0pt) -- (-2pt,0pt) node[left] {\footnotesize $\y$};
			\node[above left] at (0,0){$O$};
		\end{tikzpicture}}
	\loigiai{
		Đường thẳng bờ đi qua gốc toạ độ $O$ và điểm $(1;2)$ nên có phương trình $y=2x$ (vẽ nét liền nên kể cả bờ).\\
		Lấy điểm $M(1;0)$ thuộc phần không bị gạch, ta có $0\leq 2$ nên phần không bị gạch là miền nghiệm của bất phương trình $y\leq 2x$.
	}
\end{ex}

%%%==============Câu_12 — y = x/2 + 1 (x - 2y + 2 = 0), tô nửa dưới - phải, bờ NÉT ĐỨT=============%%%
\begin{ex}%[VNPmath - tự biên soạn]%[GV. Vũ Ngọc Phát]%[0D2H1-2]
	\immini[thm]{Phần không bị gạch (không kể đường thẳng $d$) trong hình vẽ là miền nghiệm của bất phương trình nào sau đây?
		\choice
		{\True $x-2y+2<0$}
		{$x-2y+2>0$}
		{$2x-y+2<0$}
		{$2x-y+2>0$}}
	{\begin{tikzpicture}[scale=.7,thick,>=stealth,declare function={xmin=-3;xmax=3;ymin=-1.5;ymax=3.5;f(\x)=1+(\x)/2;}]
			\begin{scope}
				\clip (xmin,ymin) rectangle (xmax,ymax);
				\fill[pattern=north east lines] (xmin,ymin) rectangle (xmax,ymax);
				\fill[white] (xmin,{f(xmin)}) -- (xmax,{f(xmax)}) -- (xmax,ymax) -- (xmin,ymax) -- cycle;
				\draw[thick,dashed] (xmin,{f(xmin)}) -- (xmax,{f(xmax)});
			\end{scope}
			\draw[->] (xmin,0) -- (xmax,0)node[above]{$x$};
			\foreach \x in {-2,-1,1,2,3}
			\draw[shift={(\x,0)}] (0pt,2pt) -- (0pt,-2pt) node[below] {\footnotesize $\x$};
			\draw[->,color=black] (0,ymin) -- (0,ymax)node[left]{$y$};
			\foreach \y in {-1,1,2,3}
			\draw[shift={(0,\y)}] (2pt,0pt) -- (-2pt,0pt) node[left] {\footnotesize $\y$};
			\node[above left] at (0,0){$O$};
		\end{tikzpicture}}
	\loigiai{
		Đường thẳng bờ $d$ đi qua hai điểm $(0;1)$ và $(2;2)$ nên có phương trình $y=\dfrac{x}{2}+1$, hay $x-2y+2=0$.\\
		$d$ vẽ nét đứt nên không kể bờ.\\
		Lấy điểm $M(0;2)$ thuộc phần không bị gạch, ta có $0-2\cdot 2+2=-2<0$ nên phần không bị gạch là miền nghiệm của bất phương trình $x-2y+2<0$.
	}
\end{ex}

%%%==============Câu_13 — y = x + 2 (x - y + 2 = 0), tô nửa trên - trái, bờ NÉT ĐỨT=============%%%
\begin{ex}%[VNPmath - tự biên soạn]%[GV. Vũ Ngọc Phát]%[0D2H1-2]
	\immini[thm]{Phần không bị gạch (không kể đường thẳng $d$) trong hình vẽ là miền nghiệm của bất phương trình nào sau đây?
		\choice
		{$x-y+2<0$}
		{\True $x-y+2>0$}
		{$x+y+2<0$}
		{$x+y+2>0$}}
	{\begin{tikzpicture}[scale=.7,thick,>=stealth,declare function={xmin=-4;xmax=2;ymin=-1.5;ymax=3.5;f(\x)=(\x)+2;}]
			\begin{scope}
				\clip (xmin,ymin) rectangle (xmax,ymax);
				\fill[pattern=north east lines] (xmin,ymin) rectangle (xmax,ymax);
				\fill[white] (xmin,{f(xmin)}) -- (xmax,{f(xmax)}) -- (xmax,ymin) -- (xmin,ymin) -- cycle;
				\draw[thick,dashed] (xmin,{f(xmin)}) -- (xmax,{f(xmax)});
			\end{scope}
			\draw[->] (xmin,0) -- (xmax,0)node[above]{$x$};
			\foreach \x in {-3,-2,-1,1,2}
			\draw[shift={(\x,0)}] (0pt,2pt) -- (0pt,-2pt) node[below] {\footnotesize $\x$};
			\draw[->,color=black] (0,ymin) -- (0,ymax)node[left]{$y$};
			\foreach \y in {-1,1,2,3}
			\draw[shift={(0,\y)}] (2pt,0pt) -- (-2pt,0pt) node[left] {\footnotesize $\y$};
			\node[above left] at (0,0){$O$};
		\end{tikzpicture}}
	\loigiai{
		Đường thẳng bờ $d$ đi qua hai điểm $(-2;0)$ và $(0;2)$ nên có phương trình $y=x+2$, hay $x-y+2=0$.\\
		$d$ vẽ nét đứt nên không kể bờ.\\
		Lấy điểm $M(0;1)$ thuộc phần không bị gạch, ta có $0-1+2=1>0$ nên phần không bị gạch là miền nghiệm của bất phương trình $x-y+2>0$.
	}
\end{ex}

%%%%%%===================================================================%%%%%%
%%%%%%============================= DẠNG 4 ===============================%%%%%%
%%%%%% Bờ là đường thẳng ax+by = c với ab > 0  (hệ số góc âm)
%%%%%%===================================================================%%%%%%
\subsubsection{Bờ là đường thẳng $ax+by=c$ với $ab>0$}

%%%==============Câu_14 — x + y = 1, tô nửa dưới - trái, bờ NÉT ĐỨT=============%%%
\begin{ex}%[VNPmath - tự biên soạn]%[GV. Vũ Ngọc Phát]%[0D2H1-2]
	\immini[thm]{Phần không bị gạch (không kể đường thẳng $d$) trong hình vẽ là miền nghiệm của bất phương trình nào sau đây?
		\choice
		{$x+y<1$}
		{$x-y>1$}
		{\True $x+y>1$}
		{$x-y<1$}}
	{\begin{tikzpicture}[scale=.7,thick,>=stealth,declare function={xmin=-2.5;xmax=3.5;ymin=-2;ymax=3;f(\x)=1-(\x);}]
			\begin{scope}
				\clip (xmin,ymin) rectangle (xmax,ymax);
				\fill[pattern=north east lines] (xmin,ymin) rectangle (xmax,ymax);
				\fill[white] (xmin,{f(xmin)}) -- (xmax,{f(xmax)}) -- (xmax,ymax) -- (xmin,ymax) -- cycle;
				\draw[thick,dashed] (xmin,{f(xmin)}) -- (xmax,{f(xmax)});
			\end{scope}
			\draw[->] (xmin,0) -- (xmax,0)node[above]{$x$};
			\foreach \x in {-2,-1,1,2,3}
			\draw[shift={(\x,0)}] (0pt,2pt) -- (0pt,-2pt) node[below] {\footnotesize $\x$};
			\draw[->,color=black] (0,ymin) -- (0,ymax)node[left]{$y$};
			\foreach \y in {-1,1,2}
			\draw[shift={(0,\y)}] (2pt,0pt) -- (-2pt,0pt) node[left] {\footnotesize $\y$};
			\node[above left] at (0,0){$O$};
		\end{tikzpicture}}
	\loigiai{
		Đường thẳng bờ $d$ đi qua hai điểm $(1;0)$ và $(0;1)$ nên có phương trình $x+y=1$.\\
		$d$ vẽ nét đứt nên không kể bờ.\\
		Lấy điểm $M(2;0)$ thuộc phần không bị gạch, ta có $2+0=2>1$ nên phần không bị gạch là miền nghiệm của bất phương trình $x+y>1$.
	}
\end{ex}

%%%==============Câu_15 — x + y = 3, tô nửa dưới - trái, bờ NÉT LIỀN=============%%%
\begin{ex}%[VNPmath - tự biên soạn]%[GV. Vũ Ngọc Phát]%[0D2H1-2]
	\immini[thm]{Phần không bị gạch trong hình vẽ là miền nghiệm của bất phương trình nào sau đây?
		\choice
		{$x+y<3$}
		{\True $x+y\geq 3$}
		{$x-y\geq 3$}
		{$x+y\leq 3$}}
	{\begin{tikzpicture}[scale=.7,thick,>=stealth,declare function={xmin=-1.5;xmax=4.5;ymin=-1;ymax=4;f(\x)=3-(\x);}]
			\begin{scope}
				\clip (xmin,ymin) rectangle (xmax,ymax);
				\fill[pattern=north east lines] (xmin,ymin) rectangle (xmax,ymax);
				\fill[white] (xmin,{f(xmin)}) -- (xmax,{f(xmax)}) -- (xmax,ymax) -- (xmin,ymax) -- cycle;
				\draw[thick] (xmin,{f(xmin)}) -- (xmax,{f(xmax)});
			\end{scope}
			\draw[->] (xmin,0) -- (xmax,0)node[above]{$x$};
			\foreach \x in {-1,1,2,3,4}
			\draw[shift={(\x,0)}] (0pt,2pt) -- (0pt,-2pt) node[below] {\footnotesize $\x$};
			\draw[->,color=black] (0,ymin) -- (0,ymax)node[left]{$y$};
			\foreach \y in {-1,1,2,3}
			\draw[shift={(0,\y)}] (2pt,0pt) -- (-2pt,0pt) node[left] {\footnotesize $\y$};
			\node[above left] at (0,0){$O$};
		\end{tikzpicture}}
	\loigiai{
		Đường thẳng bờ đi qua hai điểm $(3;0)$ và $(0;3)$ nên có phương trình $x+y=3$ (vẽ nét liền nên kể cả bờ).\\
		Lấy điểm $M(4;0)$ thuộc phần không bị gạch, ta có $4+0=4\geq 3$ nên phần không bị gạch là miền nghiệm của bất phương trình $x+y\geq 3$.
	}
\end{ex}

%%%==============Câu_16 — x + 2y = 2, tô nửa trên - phải, bờ NÉT ĐỨT=============%%%
\begin{ex}%[VNPmath - tự biên soạn]%[GV. Vũ Ngọc Phát]%[0D2H1-2]
	\immini[thm]{Phần không bị gạch (không kể đường thẳng $d$) trong hình vẽ là miền nghiệm của bất phương trình nào sau đây?
		\choice
		{$x+2y>2$}
		{$2x+y>2$}
		{\True $x+2y<2$}
		{$2x+y<2$}}
	{\begin{tikzpicture}[scale=.7,thick,>=stealth,declare function={xmin=-2;xmax=4;ymin=-2;ymax=3;f(\x)=1-(\x)/2;}]
			\begin{scope}
				\clip (xmin,ymin) rectangle (xmax,ymax);
				\fill[pattern=north east lines] (xmin,ymin) rectangle (xmax,ymax);
				\fill[white] (xmin,{f(xmin)}) -- (xmax,{f(xmax)}) -- (xmax,ymin) -- (xmin,ymin) -- cycle;
				\draw[thick,dashed] (xmin,{f(xmin)}) -- (xmax,{f(xmax)});
			\end{scope}
			\draw[->] (xmin,0) -- (xmax,0)node[above]{$x$};
			\foreach \x in {-1,1,2,3,4}
			\draw[shift={(\x,0)}] (0pt,2pt) -- (0pt,-2pt) node[below] {\footnotesize $\x$};
			\draw[->,color=black] (0,ymin) -- (0,ymax)node[left]{$y$};
			\foreach \y in {-1,1,2}
			\draw[shift={(0,\y)}] (2pt,0pt) -- (-2pt,0pt) node[left] {\footnotesize $\y$};
			\node[above left] at (0,0){$O$};
		\end{tikzpicture}}
	\loigiai{
		Đường thẳng bờ $d$ đi qua hai điểm $(0;1)$ và $(2;0)$ nên có phương trình $y=1-\dfrac{x}{2}$, hay $x+2y=2$.\\
		$d$ vẽ nét đứt nên không kể bờ.\\
		Lấy điểm $O(0;0)$ thuộc phần không bị gạch, ta có $0+2\cdot 0=0<2$ nên phần không bị gạch là miền nghiệm của bất phương trình $x+2y<2$.
	}
\end{ex}


\Closesolutionfile{ans}

% 2) Nếu muốn GỘP vào TN-0D2-B1.tex: xoá 2 dòng \Opensolutionfile/\Closesolutionfile
%    ở file này rồi dán phần các câu hỏi vào TRƯỚC dòng \Closesolutionfile{ans} của file đó.
% 3) In bảng đáp án: \inputansbox{10}{ans/ans-0D2-B1-MienNghiem}
% 4) Lời giải ngắn đã điền sẵn trong \loigiai{}; xoá nội dung nếu muốn để trống như mẫu gốc.
%    QUY ƯỚC LỜI GIẢI: chỉ phân tích miền KHÔNG bị gạch — nêu phương trình đường bờ (kèm ghi
%    chú nét liền/nét đứt), rồi lấy điểm $M$ thuộc phần không bị gạch (ưu tiên điểm có ít nhất
%    một toạ độ bằng 0) thay vào bất phương trình để kết luận.
%    - Nếu lấy gốc toạ độ thì gọi là $O(0;0)$ (KHÔNG dùng $M(0;0)$).
% 5) Đang dùng \subsubsection làm tiêu đề dạng (đánh số 1., 2., 3., 4.; KHÔNG reset số câu,
%    nên bảng đáp án vẫn liên tục 1 → 16). Nếu muốn khung "DẠNG TOÁN" thì thay bằng
%    \begin{dang}{...} ... \end{dang} — lưu ý môi trường dang tự reset số câu về 1.
% 6) Quy ước nhãn câu: %[Nguồn]%[Người biên soạn]%[Mã ID].
%    Mã ID dùng ở đây: 0D2N1-2 (nhận biết — bờ song song trục) và 0D2H1-2 (thông hiểu — bờ xiên),
%    khớp với data/10/0D2/0-KNTT-C2B3-X.tex.
%
% GHI CHÚ: hai câu bờ x = m, y = m dạng nét đứt trước đây dùng m = 3 (khung phải lệch, khó đạt
% 50% mà vẫn thấy gốc O) đã đổi thành m = -2 (câu 5) và m = -1 (câu 9) để khung đối xứng qua d.
%%%%%%%%%%%%%%%%%%%%%%%%%%%%%%%%%%%%%%%%%%%%%%%%%%%%%%%%%%%%%%%%%%%%%%%%%%%%%%%%%%%%%

\Opensolutionfile{ans}[ans/ans-0D2-B1-MienNghiem]

%%%%%%===================================================================%%%%%%
%%%%%%============================= DẠNG 1 ===============================%%%%%%
%%%%%% Bờ là đường thẳng song song với trục tung: x = m
%%%%%%===================================================================%%%%%%
\subsubsection{Bờ là đường thẳng $x=m$}

%%%==============Câu_1 — x = 0 (trục tung), tô bên trái, bờ NÉT LIỀN=============%%%
\begin{ex}%[VNPmath - tự biên soạn]%[GV. Vũ Ngọc Phát]%[0D2N1-2]
	\immini[thm]{Phần không bị gạch trong hình vẽ là miền nghiệm của bất phương trình nào sau đây?
		\choice
		{$x\leq 0$}
		{\True $x\geq 0$}
		{$y\leq 0$}
		{$y\geq 0$}}
	{\begin{tikzpicture}[scale=.7,thick,>=stealth,declare function={xmin=-3;xmax=3;ymin=-2;ymax=3;m=0;}]
			\fill[pattern=north east lines] (xmin,ymin) rectangle (xmax,ymax);
			\fill[white] (m,ymin) rectangle (xmax,ymax);
			\draw[->] (xmin,0) -- (xmax,0)node[above]{$x$};
			\foreach \x in {-2,-1,1,2,3}
			\draw[shift={(\x,0)}] (0pt,2pt) -- (0pt,-2pt) node[below] {\footnotesize $\x$};
			\draw[->,color=black] (0,ymin) -- (0,ymax)node[left]{$y$};
			\foreach \y in {-1,1,2}
			\draw[shift={(0,\y)}] (2pt,0pt) -- (-2pt,0pt) node[left] {\footnotesize $\y$};
			\node[above left] at (0,0){$O$};
		\end{tikzpicture}}
	\loigiai{
		Đường thẳng bờ là trục tung $d\colon x=0$ (vẽ nét liền nên kể cả bờ).\\
		Lấy điểm $M(1;0)$ thuộc phần không bị gạch, ta có $1\geq 0$ nên phần không bị gạch là miền nghiệm của bất phương trình $x\geq 0$.
	}
\end{ex}

%%%==============Câu_2 — x = 1, tô bên trái, bờ NÉT LIỀN=============%%%
\begin{ex}%[VNPmath - tự biên soạn]%[GV. Vũ Ngọc Phát]%[0D2N1-2]
	\immini[thm]{Phần không bị gạch trong hình vẽ là miền nghiệm của bất phương trình nào sau đây?
		\choice
		{\True $x\geq 1$}
		{$x\leq 1$}
		{$x\geq -1$}
		{$x\leq -1$}}
	{\begin{tikzpicture}[scale=.7,thick,>=stealth,declare function={xmin=-2;xmax=4;ymin=-2;ymax=3;m=1;}]
			\fill[pattern=north east lines] (xmin,ymin) rectangle (xmax,ymax);
			\fill[white] (m,ymin) rectangle (xmax,ymax);
			\draw[thick] (m,ymin) -- (m,ymax);
			\draw[->] (xmin,0) -- (xmax,0)node[above]{$x$};
			\foreach \x in {-1,1,2,3,4}
			\draw[shift={(\x,0)}] (0pt,2pt) -- (0pt,-2pt) node[below] {\footnotesize $\x$};
			\draw[->,color=black] (0,ymin) -- (0,ymax)node[left]{$y$};
			\foreach \y in {-1,1,2}
			\draw[shift={(0,\y)}] (2pt,0pt) -- (-2pt,0pt) node[left] {\footnotesize $\y$};
			\node[above left] at (0,0){$O$};
		\end{tikzpicture}}
	\loigiai{
		Đường thẳng bờ $d\colon x=1$ (vẽ nét liền nên kể cả bờ).\\
		Lấy điểm $M(2;0)$ thuộc phần không bị gạch, ta có $2\geq 1$ nên phần không bị gạch là miền nghiệm của bất phương trình $x\geq 1$.
	}
\end{ex}

%%%==============Câu_3 — x = -1, tô bên phải, bờ NÉT LIỀN=============%%%
\begin{ex}%[VNPmath - tự biên soạn]%[GV. Vũ Ngọc Phát]%[0D2N1-2]
	\immini[thm]{Phần không bị gạch trong hình vẽ là miền nghiệm của bất phương trình nào sau đây?
		\choice
		{$x\geq -1$}
		{$x\geq 1$}
		{\True $x\leq -1$}
		{$x\leq 1$}}
	{\begin{tikzpicture}[scale=.7,thick,>=stealth,declare function={xmin=-4;xmax=2;ymin=-2;ymax=3;m=-1;}]
			\fill[pattern=north east lines] (xmin,ymin) rectangle (xmax,ymax);
			\fill[white] (xmin,ymin) rectangle (m,ymax);
			\draw[thick] (m,ymin) -- (m,ymax);
			\draw[->] (xmin,0) -- (xmax,0)node[above]{$x$};
			\foreach \x in {-3,-2,-1,1,2}
			\draw[shift={(\x,0)}] (0pt,2pt) -- (0pt,-2pt) node[below] {\footnotesize $\x$};
			\draw[->,color=black] (0,ymin) -- (0,ymax)node[left]{$y$};
			\foreach \y in {-1,1,2}
			\draw[shift={(0,\y)}] (2pt,0pt) -- (-2pt,0pt) node[left] {\footnotesize $\y$};
			\node[above left] at (0,0){$O$};
		\end{tikzpicture}}
	\loigiai{
		Đường thẳng bờ $d\colon x=-1$ (vẽ nét liền nên kể cả bờ).\\
		Lấy điểm $M(-2;0)$ thuộc phần không bị gạch, ta có $-2\leq -1$ nên phần không bị gạch là miền nghiệm của bất phương trình $x\leq -1$.
	}
\end{ex}

%%%==============Câu_4 — x = 2, tô bên phải, bờ NÉT ĐỨT=============%%%
\begin{ex}%[VNPmath - tự biên soạn]%[GV. Vũ Ngọc Phát]%[0D2N1-2]
	\immini[thm]{Phần không bị gạch (không kể đường thẳng $d$) trong hình vẽ là miền nghiệm của bất phương trình nào sau đây?
		\choice
		{$x>2$}
		{$x<3$}
		{$x>3$}
		{\True $x<2$}}
	{\begin{tikzpicture}[scale=.7,thick,>=stealth,declare function={xmin=-1;xmax=5;ymin=-2;ymax=3;m=2;}]
			\fill[pattern=north east lines] (xmin,ymin) rectangle (xmax,ymax);
			\fill[white] (xmin,ymin) rectangle (m,ymax);
			\draw[thick,dashed] (m,ymin) -- (m,ymax);
			\draw[->] (xmin,0) -- (xmax,0)node[above]{$x$};
			\foreach \x in {-1,1,2,3,4,5}
			\draw[shift={(\x,0)}] (0pt,2pt) -- (0pt,-2pt) node[below] {\footnotesize $\x$};
			\draw[->,color=black] (0,ymin) -- (0,ymax)node[left]{$y$};
			\foreach \y in {-1,1,2}
			\draw[shift={(0,\y)}] (2pt,0pt) -- (-2pt,0pt) node[left] {\footnotesize $\y$};
			\node[above left] at (0,0){$O$};
		\end{tikzpicture}}
	\loigiai{
		Đường thẳng bờ $d\colon x=2$ vẽ nét đứt nên không kể bờ.\\
		Lấy điểm $O(0;0)$ thuộc phần không bị gạch, ta có $0<2$ nên phần không bị gạch là miền nghiệm của bất phương trình $x<2$.
	}
\end{ex}

%%%==============Câu_5 — x = -2, tô bên trái, bờ NÉT ĐỨT=============%%%
\begin{ex}%[VNPmath - tự biên soạn]%[GV. Vũ Ngọc Phát]%[0D2N1-2]
	\immini[thm]{Phần không bị gạch (không kể đường thẳng $d$) trong hình vẽ là miền nghiệm của bất phương trình nào sau đây?
		\choice
		{$x<-3$}
		{$x>-3$}
		{$x<-2$}
		{\True $x>-2$}}
	{\begin{tikzpicture}[scale=.7,thick,>=stealth,declare function={xmin=-5;xmax=1;ymin=-2;ymax=3;m=-2;}]
			\fill[pattern=north east lines] (xmin,ymin) rectangle (xmax,ymax);
			\fill[white] (m,ymin) rectangle (xmax,ymax);
			\draw[thick,dashed] (m,ymin) -- (m,ymax);
			\draw[->] (xmin,0) -- (xmax,0)node[above]{$x$};
			\foreach \x in {-4,-3,-2,-1,1}
			\draw[shift={(\x,0)}] (0pt,2pt) -- (0pt,-2pt) node[below] {\footnotesize $\x$};
			\draw[->,color=black] (0,ymin) -- (0,ymax)node[left]{$y$};
			\foreach \y in {-1,1,2}
			\draw[shift={(0,\y)}] (2pt,0pt) -- (-2pt,0pt) node[left] {\footnotesize $\y$};
			\node[above left] at (0,0){$O$};
		\end{tikzpicture}}
	\loigiai{
		Đường thẳng bờ $d\colon x=-2$ vẽ nét đứt nên không kể bờ.\\
		Lấy điểm $O(0;0)$ thuộc phần không bị gạch, ta có $0>-2$ nên phần không bị gạch là miền nghiệm của bất phương trình $x>-2$.
	}
\end{ex}

%%%%%%===================================================================%%%%%%
%%%%%%============================= DẠNG 2 ===============================%%%%%%
%%%%%% Bờ là đường thẳng song song với trục hoành: y = m
%%%%%%===================================================================%%%%%%
\subsubsection{Bờ là đường thẳng $y=m$}

%%%==============Câu_6 — y = 0 (trục hoành), tô phía dưới, bờ NÉT LIỀN=============%%%
\begin{ex}%[VNPmath - tự biên soạn]%[GV. Vũ Ngọc Phát]%[0D2N1-2]
	\immini[thm]{Phần không bị gạch trong hình vẽ là miền nghiệm của bất phương trình nào sau đây?
		\choice
		{$y\leq 0$}
		{$x\leq 0$}
		{$x\geq 0$}
		{\True $y\geq 0$}}
	{\begin{tikzpicture}[scale=.7,thick,>=stealth,declare function={xmin=-2;xmax=4;ymin=-2.5;ymax=2.5;m=0;}]
			\fill[pattern=north east lines] (xmin,ymin) rectangle (xmax,ymax);
			\fill[white] (xmin,m) rectangle (xmax,ymax);
			\draw[->] (xmin,0) -- (xmax,0)node[above]{$x$};
			\foreach \x in {-1,1,2,3,4}
			\draw[shift={(\x,0)}] (0pt,2pt) -- (0pt,-2pt) node[below] {\footnotesize $\x$};
			\draw[->,color=black] (0,ymin) -- (0,ymax)node[left]{$y$};
			\foreach \y in {-2,-1,1,2}
			\draw[shift={(0,\y)}] (2pt,0pt) -- (-2pt,0pt) node[left] {\footnotesize $\y$};
			\node[above left] at (0,0){$O$};
		\end{tikzpicture}}
	\loigiai{
		Đường thẳng bờ là trục hoành $d\colon y=0$ (vẽ nét liền nên kể cả bờ).\\
		Lấy điểm $M(0;1)$ thuộc phần không bị gạch, ta có $1\geq 0$ nên phần không bị gạch là miền nghiệm của bất phương trình $y\geq 0$.
	}
\end{ex}

%%%==============Câu_7 — y = 1, tô phía dưới, bờ NÉT ĐỨT=============%%%
\begin{ex}%[VNPmath - tự biên soạn]%[GV. Vũ Ngọc Phát]%[0D2N1-2]
	\immini[thm]{Phần không bị gạch (không kể đường thẳng $d$) trong hình vẽ là miền nghiệm của bất phương trình nào sau đây?
		\choice
		{\True $y>1$}
		{$y<1$}
		{$y>2$}
		{$y<2$}}
	{\begin{tikzpicture}[scale=.7,thick,>=stealth,declare function={xmin=-2;xmax=4;ymin=-1.5;ymax=3.5;m=1;}]
			\fill[pattern=north east lines] (xmin,ymin) rectangle (xmax,ymax);
			\fill[white] (xmin,m) rectangle (xmax,ymax);
			\draw[thick,dashed] (xmin,m) -- (xmax,m);
			\draw[->] (xmin,0) -- (xmax,0)node[above]{$x$};
			\foreach \x in {-1,1,2,3,4}
			\draw[shift={(\x,0)}] (0pt,2pt) -- (0pt,-2pt) node[below] {\footnotesize $\x$};
			\draw[->,color=black] (0,ymin) -- (0,ymax)node[left]{$y$};
			\foreach \y in {-1,1,2,3}
			\draw[shift={(0,\y)}] (2pt,0pt) -- (-2pt,0pt) node[left] {\footnotesize $\y$};
			\node[above left] at (0,0){$O$};
		\end{tikzpicture}}
	\loigiai{
		Đường thẳng bờ $d\colon y=1$ vẽ nét đứt nên không kể bờ.\\
		Lấy điểm $M(0;2)$ thuộc phần không bị gạch, ta có $2>1$ nên phần không bị gạch là miền nghiệm của bất phương trình $y>1$.
	}
\end{ex}

%%%==============Câu_8 — y = 2, tô phía trên, bờ NÉT LIỀN=============%%%
\begin{ex}%[VNPmath - tự biên soạn]%[GV. Vũ Ngọc Phát]%[0D2N1-2]
	\immini[thm]{Phần không bị gạch trong hình vẽ là miền nghiệm của bất phương trình nào sau đây?
		\choice
		{$y\geq 2$}
		{$y\geq -2$}
		{\True $y\leq 2$}
		{$y\leq -2$}}
	{\begin{tikzpicture}[scale=.7,thick,>=stealth,declare function={xmin=-2;xmax=4;ymin=-0.5;ymax=4.5;m=2;}]
			\fill[pattern=north east lines] (xmin,ymin) rectangle (xmax,ymax);
			\fill[white] (xmin,ymin) rectangle (xmax,m);
			\draw[thick] (xmin,m) -- (xmax,m);
			\draw[->] (xmin,0) -- (xmax,0)node[above]{$x$};
			\foreach \x in {-1,1,2,3,4}
			\draw[shift={(\x,0)}] (0pt,2pt) -- (0pt,-2pt) node[below] {\footnotesize $\x$};
			\draw[->,color=black] (0,ymin) -- (0,ymax)node[left]{$y$};
			\foreach \y in {1,2,3,4}
			\draw[shift={(0,\y)}] (2pt,0pt) -- (-2pt,0pt) node[left] {\footnotesize $\y$};
			\node[above left] at (0,0){$O$};
		\end{tikzpicture}}
	\loigiai{
		Đường thẳng bờ $d\colon y=2$ (vẽ nét liền nên kể cả bờ).\\
		Lấy điểm $O(0;0)$ thuộc phần không bị gạch, ta có $0\leq 2$ nên phần không bị gạch là miền nghiệm của bất phương trình $y\leq 2$.
	}
\end{ex}

%%%==============Câu_9 — y = -1, tô phía trên, bờ NÉT ĐỨT=============%%%
\begin{ex}%[VNPmath - tự biên soạn]%[GV. Vũ Ngọc Phát]%[0D2N1-2]
	\immini[thm]{Phần không bị gạch (không kể đường thẳng $d$) trong hình vẽ là miền nghiệm của bất phương trình nào sau đây?
		\choice
		{$y>-1$}
		{\True $y<-1$}
		{$y>-2$}
		{$y<-2$}}
	{\begin{tikzpicture}[scale=.7,thick,>=stealth,declare function={xmin=-2;xmax=4;ymin=-3.5;ymax=1.5;m=-1;}]
			\fill[pattern=north east lines] (xmin,ymin) rectangle (xmax,ymax);
			\fill[white] (xmin,ymin) rectangle (xmax,m);
			\draw[thick,dashed] (xmin,m) -- (xmax,m);
			\draw[->] (xmin,0) -- (xmax,0)node[above]{$x$};
			\foreach \x in {-1,1,2,3,4}
			\draw[shift={(\x,0)}] (0pt,2pt) -- (0pt,-2pt) node[below] {\footnotesize $\x$};
			\draw[->,color=black] (0,ymin) -- (0,ymax)node[left]{$y$};
			\foreach \y in {-3,-2,-1,1}
			\draw[shift={(0,\y)}] (2pt,0pt) -- (-2pt,0pt) node[left] {\footnotesize $\y$};
			\node[above left] at (0,0){$O$};
		\end{tikzpicture}}
	\loigiai{
		Đường thẳng bờ $d\colon y=-1$ vẽ nét đứt nên không kể bờ.\\
		Lấy điểm $M(0;-2)$ thuộc phần không bị gạch, ta có $-2<-1$ nên phần không bị gạch là miền nghiệm của bất phương trình $y<-1$.
	}
\end{ex}

%%%%%%===================================================================%%%%%%
%%%%%%============================= DẠNG 3 ===============================%%%%%%
%%%%%% Bờ là đường thẳng ax+by = c với ab < 0  (hệ số góc dương)
%%%%%%===================================================================%%%%%%
\subsubsection{Bờ là đường thẳng $ax+by=c$ với $ab<0$}

%%%==============Câu_10 — y = x (x - y = 0), tô nửa trên - trái, bờ NÉT LIỀN=============%%%
\begin{ex}%[VNPmath - tự biên soạn]%[GV. Vũ Ngọc Phát]%[0D2H1-2]
	\immini[thm]{Phần không bị gạch trong hình vẽ là miền nghiệm của bất phương trình nào sau đây?
		\choice
		{$y\geq x$}
		{$y\geq -x$}
		{$y\leq -x$}
		{\True $y\leq x$}}
	{\begin{tikzpicture}[scale=.7,thick,>=stealth,declare function={xmin=-2.5;xmax=3.5;ymin=-2;ymax=3;f(\x)=\x;}]
			\begin{scope}
				\clip (xmin,ymin) rectangle (xmax,ymax);
				\fill[pattern=north east lines] (xmin,ymin) rectangle (xmax,ymax);
				\fill[white] (xmin,{f(xmin)}) -- (xmax,{f(xmax)}) -- (xmax,ymin) -- (xmin,ymin) -- cycle;
				\draw[thick] (xmin,{f(xmin)}) -- (xmax,{f(xmax)});
			\end{scope}
			\draw[->] (xmin,0) -- (xmax,0)node[above]{$x$};
			\foreach \x in {-2,-1,1,2,3}
			\draw[shift={(\x,0)}] (0pt,2pt) -- (0pt,-2pt) node[below] {\footnotesize $\x$};
			\draw[->,color=black] (0,ymin) -- (0,ymax)node[left]{$y$};
			\foreach \y in {-1,1,2}
			\draw[shift={(0,\y)}] (2pt,0pt) -- (-2pt,0pt) node[left] {\footnotesize $\y$};
			\node[above left] at (0,0){$O$};
		\end{tikzpicture}}
	\loigiai{
		Đường thẳng bờ đi qua gốc toạ độ $O$ và điểm $(1;1)$ nên có phương trình $y=x$ (vẽ nét liền nên kể cả bờ).\\
		Lấy điểm $M(1;0)$ thuộc phần không bị gạch, ta có $0\leq 1$ nên phần không bị gạch là miền nghiệm của bất phương trình $y\leq x$.
	}
\end{ex}

%%%==============Câu_11 — y = 2x (2x - y = 0), tô nửa trên - trái, bờ NÉT LIỀN=============%%%
\begin{ex}%[VNPmath - tự biên soạn]%[GV. Vũ Ngọc Phát]%[0D2H1-2]
	\immini[thm]{Phần không bị gạch trong hình vẽ là miền nghiệm của bất phương trình nào sau đây?
		\choice
		{$y\geq 2x$}
		{$2y\geq x$}
		{$2y\leq x$}
		{\True $y\leq 2x$}}
	{\begin{tikzpicture}[scale=.7,thick,>=stealth,declare function={xmin=-2.5;xmax=3.5;ymin=-1.5;ymax=3.5;f(\x)=2*(\x);}]
			\begin{scope}
				\clip (xmin,ymin) rectangle (xmax,ymax);
				\fill[pattern=north east lines] (xmin,ymin) rectangle (xmax,ymax);
				\fill[white] (xmin,{f(xmin)}) -- (xmax,{f(xmax)}) -- (xmax,ymin) -- (xmin,ymin) -- cycle;
				\draw[thick] (xmin,{f(xmin)}) -- (xmax,{f(xmax)});
			\end{scope}
			\draw[->] (xmin,0) -- (xmax,0)node[above]{$x$};
			\foreach \x in {-2,-1,1,2,3}
			\draw[shift={(\x,0)}] (0pt,2pt) -- (0pt,-2pt) node[below] {\footnotesize $\x$};
			\draw[->,color=black] (0,ymin) -- (0,ymax)node[left]{$y$};
			\foreach \y in {-1,1,2,3}
			\draw[shift={(0,\y)}] (2pt,0pt) -- (-2pt,0pt) node[left] {\footnotesize $\y$};
			\node[above left] at (0,0){$O$};
		\end{tikzpicture}}
	\loigiai{
		Đường thẳng bờ đi qua gốc toạ độ $O$ và điểm $(1;2)$ nên có phương trình $y=2x$ (vẽ nét liền nên kể cả bờ).\\
		Lấy điểm $M(1;0)$ thuộc phần không bị gạch, ta có $0\leq 2$ nên phần không bị gạch là miền nghiệm của bất phương trình $y\leq 2x$.
	}
\end{ex}

%%%==============Câu_12 — y = x/2 + 1 (x - 2y + 2 = 0), tô nửa dưới - phải, bờ NÉT ĐỨT=============%%%
\begin{ex}%[VNPmath - tự biên soạn]%[GV. Vũ Ngọc Phát]%[0D2H1-2]
	\immini[thm]{Phần không bị gạch (không kể đường thẳng $d$) trong hình vẽ là miền nghiệm của bất phương trình nào sau đây?
		\choice
		{\True $x-2y+2<0$}
		{$x-2y+2>0$}
		{$2x-y+2<0$}
		{$2x-y+2>0$}}
	{\begin{tikzpicture}[scale=.7,thick,>=stealth,declare function={xmin=-3;xmax=3;ymin=-1.5;ymax=3.5;f(\x)=1+(\x)/2;}]
			\begin{scope}
				\clip (xmin,ymin) rectangle (xmax,ymax);
				\fill[pattern=north east lines] (xmin,ymin) rectangle (xmax,ymax);
				\fill[white] (xmin,{f(xmin)}) -- (xmax,{f(xmax)}) -- (xmax,ymax) -- (xmin,ymax) -- cycle;
				\draw[thick,dashed] (xmin,{f(xmin)}) -- (xmax,{f(xmax)});
			\end{scope}
			\draw[->] (xmin,0) -- (xmax,0)node[above]{$x$};
			\foreach \x in {-2,-1,1,2,3}
			\draw[shift={(\x,0)}] (0pt,2pt) -- (0pt,-2pt) node[below] {\footnotesize $\x$};
			\draw[->,color=black] (0,ymin) -- (0,ymax)node[left]{$y$};
			\foreach \y in {-1,1,2,3}
			\draw[shift={(0,\y)}] (2pt,0pt) -- (-2pt,0pt) node[left] {\footnotesize $\y$};
			\node[above left] at (0,0){$O$};
		\end{tikzpicture}}
	\loigiai{
		Đường thẳng bờ $d$ đi qua hai điểm $(0;1)$ và $(2;2)$ nên có phương trình $y=\dfrac{x}{2}+1$, hay $x-2y+2=0$.\\
		$d$ vẽ nét đứt nên không kể bờ.\\
		Lấy điểm $M(0;2)$ thuộc phần không bị gạch, ta có $0-2\cdot 2+2=-2<0$ nên phần không bị gạch là miền nghiệm của bất phương trình $x-2y+2<0$.
	}
\end{ex}

%%%==============Câu_13 — y = x + 2 (x - y + 2 = 0), tô nửa trên - trái, bờ NÉT ĐỨT=============%%%
\begin{ex}%[VNPmath - tự biên soạn]%[GV. Vũ Ngọc Phát]%[0D2H1-2]
	\immini[thm]{Phần không bị gạch (không kể đường thẳng $d$) trong hình vẽ là miền nghiệm của bất phương trình nào sau đây?
		\choice
		{$x-y+2<0$}
		{\True $x-y+2>0$}
		{$x+y+2<0$}
		{$x+y+2>0$}}
	{\begin{tikzpicture}[scale=.7,thick,>=stealth,declare function={xmin=-4;xmax=2;ymin=-1.5;ymax=3.5;f(\x)=(\x)+2;}]
			\begin{scope}
				\clip (xmin,ymin) rectangle (xmax,ymax);
				\fill[pattern=north east lines] (xmin,ymin) rectangle (xmax,ymax);
				\fill[white] (xmin,{f(xmin)}) -- (xmax,{f(xmax)}) -- (xmax,ymin) -- (xmin,ymin) -- cycle;
				\draw[thick,dashed] (xmin,{f(xmin)}) -- (xmax,{f(xmax)});
			\end{scope}
			\draw[->] (xmin,0) -- (xmax,0)node[above]{$x$};
			\foreach \x in {-3,-2,-1,1,2}
			\draw[shift={(\x,0)}] (0pt,2pt) -- (0pt,-2pt) node[below] {\footnotesize $\x$};
			\draw[->,color=black] (0,ymin) -- (0,ymax)node[left]{$y$};
			\foreach \y in {-1,1,2,3}
			\draw[shift={(0,\y)}] (2pt,0pt) -- (-2pt,0pt) node[left] {\footnotesize $\y$};
			\node[above left] at (0,0){$O$};
		\end{tikzpicture}}
	\loigiai{
		Đường thẳng bờ $d$ đi qua hai điểm $(-2;0)$ và $(0;2)$ nên có phương trình $y=x+2$, hay $x-y+2=0$.\\
		$d$ vẽ nét đứt nên không kể bờ.\\
		Lấy điểm $M(0;1)$ thuộc phần không bị gạch, ta có $0-1+2=1>0$ nên phần không bị gạch là miền nghiệm của bất phương trình $x-y+2>0$.
	}
\end{ex}

%%%%%%===================================================================%%%%%%
%%%%%%============================= DẠNG 4 ===============================%%%%%%
%%%%%% Bờ là đường thẳng ax+by = c với ab > 0  (hệ số góc âm)
%%%%%%===================================================================%%%%%%
\subsubsection{Bờ là đường thẳng $ax+by=c$ với $ab>0$}

%%%==============Câu_14 — x + y = 1, tô nửa dưới - trái, bờ NÉT ĐỨT=============%%%
\begin{ex}%[VNPmath - tự biên soạn]%[GV. Vũ Ngọc Phát]%[0D2H1-2]
	\immini[thm]{Phần không bị gạch (không kể đường thẳng $d$) trong hình vẽ là miền nghiệm của bất phương trình nào sau đây?
		\choice
		{$x+y<1$}
		{$x-y>1$}
		{\True $x+y>1$}
		{$x-y<1$}}
	{\begin{tikzpicture}[scale=.7,thick,>=stealth,declare function={xmin=-2.5;xmax=3.5;ymin=-2;ymax=3;f(\x)=1-(\x);}]
			\begin{scope}
				\clip (xmin,ymin) rectangle (xmax,ymax);
				\fill[pattern=north east lines] (xmin,ymin) rectangle (xmax,ymax);
				\fill[white] (xmin,{f(xmin)}) -- (xmax,{f(xmax)}) -- (xmax,ymax) -- (xmin,ymax) -- cycle;
				\draw[thick,dashed] (xmin,{f(xmin)}) -- (xmax,{f(xmax)});
			\end{scope}
			\draw[->] (xmin,0) -- (xmax,0)node[above]{$x$};
			\foreach \x in {-2,-1,1,2,3}
			\draw[shift={(\x,0)}] (0pt,2pt) -- (0pt,-2pt) node[below] {\footnotesize $\x$};
			\draw[->,color=black] (0,ymin) -- (0,ymax)node[left]{$y$};
			\foreach \y in {-1,1,2}
			\draw[shift={(0,\y)}] (2pt,0pt) -- (-2pt,0pt) node[left] {\footnotesize $\y$};
			\node[above left] at (0,0){$O$};
		\end{tikzpicture}}
	\loigiai{
		Đường thẳng bờ $d$ đi qua hai điểm $(1;0)$ và $(0;1)$ nên có phương trình $x+y=1$.\\
		$d$ vẽ nét đứt nên không kể bờ.\\
		Lấy điểm $M(2;0)$ thuộc phần không bị gạch, ta có $2+0=2>1$ nên phần không bị gạch là miền nghiệm của bất phương trình $x+y>1$.
	}
\end{ex}

%%%==============Câu_15 — x + y = 3, tô nửa dưới - trái, bờ NÉT LIỀN=============%%%
\begin{ex}%[VNPmath - tự biên soạn]%[GV. Vũ Ngọc Phát]%[0D2H1-2]
	\immini[thm]{Phần không bị gạch trong hình vẽ là miền nghiệm của bất phương trình nào sau đây?
		\choice
		{$x+y<3$}
		{\True $x+y\geq 3$}
		{$x-y\geq 3$}
		{$x+y\leq 3$}}
	{\begin{tikzpicture}[scale=.7,thick,>=stealth,declare function={xmin=-1.5;xmax=4.5;ymin=-1;ymax=4;f(\x)=3-(\x);}]
			\begin{scope}
				\clip (xmin,ymin) rectangle (xmax,ymax);
				\fill[pattern=north east lines] (xmin,ymin) rectangle (xmax,ymax);
				\fill[white] (xmin,{f(xmin)}) -- (xmax,{f(xmax)}) -- (xmax,ymax) -- (xmin,ymax) -- cycle;
				\draw[thick] (xmin,{f(xmin)}) -- (xmax,{f(xmax)});
			\end{scope}
			\draw[->] (xmin,0) -- (xmax,0)node[above]{$x$};
			\foreach \x in {-1,1,2,3,4}
			\draw[shift={(\x,0)}] (0pt,2pt) -- (0pt,-2pt) node[below] {\footnotesize $\x$};
			\draw[->,color=black] (0,ymin) -- (0,ymax)node[left]{$y$};
			\foreach \y in {-1,1,2,3}
			\draw[shift={(0,\y)}] (2pt,0pt) -- (-2pt,0pt) node[left] {\footnotesize $\y$};
			\node[above left] at (0,0){$O$};
		\end{tikzpicture}}
	\loigiai{
		Đường thẳng bờ đi qua hai điểm $(3;0)$ và $(0;3)$ nên có phương trình $x+y=3$ (vẽ nét liền nên kể cả bờ).\\
		Lấy điểm $M(4;0)$ thuộc phần không bị gạch, ta có $4+0=4\geq 3$ nên phần không bị gạch là miền nghiệm của bất phương trình $x+y\geq 3$.
	}
\end{ex}

%%%==============Câu_16 — x + 2y = 2, tô nửa trên - phải, bờ NÉT ĐỨT=============%%%
\begin{ex}%[VNPmath - tự biên soạn]%[GV. Vũ Ngọc Phát]%[0D2H1-2]
	\immini[thm]{Phần không bị gạch (không kể đường thẳng $d$) trong hình vẽ là miền nghiệm của bất phương trình nào sau đây?
		\choice
		{$x+2y>2$}
		{$2x+y>2$}
		{\True $x+2y<2$}
		{$2x+y<2$}}
	{\begin{tikzpicture}[scale=.7,thick,>=stealth,declare function={xmin=-2;xmax=4;ymin=-2;ymax=3;f(\x)=1-(\x)/2;}]
			\begin{scope}
				\clip (xmin,ymin) rectangle (xmax,ymax);
				\fill[pattern=north east lines] (xmin,ymin) rectangle (xmax,ymax);
				\fill[white] (xmin,{f(xmin)}) -- (xmax,{f(xmax)}) -- (xmax,ymin) -- (xmin,ymin) -- cycle;
				\draw[thick,dashed] (xmin,{f(xmin)}) -- (xmax,{f(xmax)});
			\end{scope}
			\draw[->] (xmin,0) -- (xmax,0)node[above]{$x$};
			\foreach \x in {-1,1,2,3,4}
			\draw[shift={(\x,0)}] (0pt,2pt) -- (0pt,-2pt) node[below] {\footnotesize $\x$};
			\draw[->,color=black] (0,ymin) -- (0,ymax)node[left]{$y$};
			\foreach \y in {-1,1,2}
			\draw[shift={(0,\y)}] (2pt,0pt) -- (-2pt,0pt) node[left] {\footnotesize $\y$};
			\node[above left] at (0,0){$O$};
		\end{tikzpicture}}
	\loigiai{
		Đường thẳng bờ $d$ đi qua hai điểm $(0;1)$ và $(2;0)$ nên có phương trình $y=1-\dfrac{x}{2}$, hay $x+2y=2$.\\
		$d$ vẽ nét đứt nên không kể bờ.\\
		Lấy điểm $O(0;0)$ thuộc phần không bị gạch, ta có $0+2\cdot 0=0<2$ nên phần không bị gạch là miền nghiệm của bất phương trình $x+2y<2$.
	}
\end{ex}


\Closesolutionfile{ans}

% 2) Nếu muốn GỘP vào TN-0D2-B1.tex: xoá 2 dòng \Opensolutionfile/\Closesolutionfile
%    ở file này rồi dán phần các câu hỏi vào TRƯỚC dòng \Closesolutionfile{ans} của file đó.
% 3) In bảng đáp án: \inputansbox{10}{ans/ans-0D2-B1-MienNghiem}
% 4) Lời giải ngắn đã điền sẵn trong \loigiai{}; xoá nội dung nếu muốn để trống như mẫu gốc.
%    QUY ƯỚC LỜI GIẢI: chỉ phân tích miền KHÔNG bị gạch — nêu phương trình đường bờ (kèm ghi
%    chú nét liền/nét đứt), rồi lấy điểm $M$ thuộc phần không bị gạch (ưu tiên điểm có ít nhất
%    một toạ độ bằng 0) thay vào bất phương trình để kết luận.
%    - Nếu lấy gốc toạ độ thì gọi là $O(0;0)$ (KHÔNG dùng $M(0;0)$).
% 5) Đang dùng \subsubsection làm tiêu đề dạng (đánh số 1., 2., 3., 4.; KHÔNG reset số câu,
%    nên bảng đáp án vẫn liên tục 1 → 16). Nếu muốn khung "DẠNG TOÁN" thì thay bằng
%    \begin{dang}{...} ... \end{dang} — lưu ý môi trường dang tự reset số câu về 1.
% 6) Quy ước nhãn câu: %[Nguồn]%[Người biên soạn]%[Mã ID].
%    Mã ID dùng ở đây: 0D2N1-2 (nhận biết — bờ song song trục) và 0D2H1-2 (thông hiểu — bờ xiên),
%    khớp với data/10/0D2/0-KNTT-C2B3-X.tex.
%
% GHI CHÚ: hai câu bờ x = m, y = m dạng nét đứt trước đây dùng m = 3 (khung phải lệch, khó đạt
% 50% mà vẫn thấy gốc O) đã đổi thành m = -2 (câu 5) và m = -1 (câu 9) để khung đối xứng qua d.
%%%%%%%%%%%%%%%%%%%%%%%%%%%%%%%%%%%%%%%%%%%%%%%%%%%%%%%%%%%%%%%%%%%%%%%%%%%%%%%%%%%%%

\Opensolutionfile{ans}[ans/ans-0D2-B1-MienNghiem]

%%%%%%===================================================================%%%%%%
%%%%%%============================= DẠNG 1 ===============================%%%%%%
%%%%%% Bờ là đường thẳng song song với trục tung: x = m
%%%%%%===================================================================%%%%%%
\subsubsection{Bờ là đường thẳng $x=m$}

%%%==============Câu_1 — x = 0 (trục tung), tô bên trái, bờ NÉT LIỀN=============%%%
\begin{ex}%[VNPmath - tự biên soạn]%[GV. Vũ Ngọc Phát]%[0D2N1-2]
	\immini[thm]{Phần không bị gạch trong hình vẽ là miền nghiệm của bất phương trình nào sau đây?
		\choice
		{$x\leq 0$}
		{\True $x\geq 0$}
		{$y\leq 0$}
		{$y\geq 0$}}
	{\begin{tikzpicture}[scale=.7,thick,>=stealth,declare function={xmin=-3;xmax=3;ymin=-2;ymax=3;m=0;}]
			\fill[pattern=north east lines] (xmin,ymin) rectangle (xmax,ymax);
			\fill[white] (m,ymin) rectangle (xmax,ymax);
			\draw[->] (xmin,0) -- (xmax,0)node[above]{$x$};
			\foreach \x in {-2,-1,1,2,3}
			\draw[shift={(\x,0)}] (0pt,2pt) -- (0pt,-2pt) node[below] {\footnotesize $\x$};
			\draw[->,color=black] (0,ymin) -- (0,ymax)node[left]{$y$};
			\foreach \y in {-1,1,2}
			\draw[shift={(0,\y)}] (2pt,0pt) -- (-2pt,0pt) node[left] {\footnotesize $\y$};
			\node[above left] at (0,0){$O$};
		\end{tikzpicture}}
	\loigiai{
		Đường thẳng bờ là trục tung $d\colon x=0$ (vẽ nét liền nên kể cả bờ).\\
		Lấy điểm $M(1;0)$ thuộc phần không bị gạch, ta có $1\geq 0$ nên phần không bị gạch là miền nghiệm của bất phương trình $x\geq 0$.
	}
\end{ex}

%%%==============Câu_2 — x = 1, tô bên trái, bờ NÉT LIỀN=============%%%
\begin{ex}%[VNPmath - tự biên soạn]%[GV. Vũ Ngọc Phát]%[0D2N1-2]
	\immini[thm]{Phần không bị gạch trong hình vẽ là miền nghiệm của bất phương trình nào sau đây?
		\choice
		{\True $x\geq 1$}
		{$x\leq 1$}
		{$x\geq -1$}
		{$x\leq -1$}}
	{\begin{tikzpicture}[scale=.7,thick,>=stealth,declare function={xmin=-2;xmax=4;ymin=-2;ymax=3;m=1;}]
			\fill[pattern=north east lines] (xmin,ymin) rectangle (xmax,ymax);
			\fill[white] (m,ymin) rectangle (xmax,ymax);
			\draw[thick] (m,ymin) -- (m,ymax);
			\draw[->] (xmin,0) -- (xmax,0)node[above]{$x$};
			\foreach \x in {-1,1,2,3,4}
			\draw[shift={(\x,0)}] (0pt,2pt) -- (0pt,-2pt) node[below] {\footnotesize $\x$};
			\draw[->,color=black] (0,ymin) -- (0,ymax)node[left]{$y$};
			\foreach \y in {-1,1,2}
			\draw[shift={(0,\y)}] (2pt,0pt) -- (-2pt,0pt) node[left] {\footnotesize $\y$};
			\node[above left] at (0,0){$O$};
		\end{tikzpicture}}
	\loigiai{
		Đường thẳng bờ $d\colon x=1$ (vẽ nét liền nên kể cả bờ).\\
		Lấy điểm $M(2;0)$ thuộc phần không bị gạch, ta có $2\geq 1$ nên phần không bị gạch là miền nghiệm của bất phương trình $x\geq 1$.
	}
\end{ex}

%%%==============Câu_3 — x = -1, tô bên phải, bờ NÉT LIỀN=============%%%
\begin{ex}%[VNPmath - tự biên soạn]%[GV. Vũ Ngọc Phát]%[0D2N1-2]
	\immini[thm]{Phần không bị gạch trong hình vẽ là miền nghiệm của bất phương trình nào sau đây?
		\choice
		{$x\geq -1$}
		{$x\geq 1$}
		{\True $x\leq -1$}
		{$x\leq 1$}}
	{\begin{tikzpicture}[scale=.7,thick,>=stealth,declare function={xmin=-4;xmax=2;ymin=-2;ymax=3;m=-1;}]
			\fill[pattern=north east lines] (xmin,ymin) rectangle (xmax,ymax);
			\fill[white] (xmin,ymin) rectangle (m,ymax);
			\draw[thick] (m,ymin) -- (m,ymax);
			\draw[->] (xmin,0) -- (xmax,0)node[above]{$x$};
			\foreach \x in {-3,-2,-1,1,2}
			\draw[shift={(\x,0)}] (0pt,2pt) -- (0pt,-2pt) node[below] {\footnotesize $\x$};
			\draw[->,color=black] (0,ymin) -- (0,ymax)node[left]{$y$};
			\foreach \y in {-1,1,2}
			\draw[shift={(0,\y)}] (2pt,0pt) -- (-2pt,0pt) node[left] {\footnotesize $\y$};
			\node[above left] at (0,0){$O$};
		\end{tikzpicture}}
	\loigiai{
		Đường thẳng bờ $d\colon x=-1$ (vẽ nét liền nên kể cả bờ).\\
		Lấy điểm $M(-2;0)$ thuộc phần không bị gạch, ta có $-2\leq -1$ nên phần không bị gạch là miền nghiệm của bất phương trình $x\leq -1$.
	}
\end{ex}

%%%==============Câu_4 — x = 2, tô bên phải, bờ NÉT ĐỨT=============%%%
\begin{ex}%[VNPmath - tự biên soạn]%[GV. Vũ Ngọc Phát]%[0D2N1-2]
	\immini[thm]{Phần không bị gạch (không kể đường thẳng $d$) trong hình vẽ là miền nghiệm của bất phương trình nào sau đây?
		\choice
		{$x>2$}
		{$x<3$}
		{$x>3$}
		{\True $x<2$}}
	{\begin{tikzpicture}[scale=.7,thick,>=stealth,declare function={xmin=-1;xmax=5;ymin=-2;ymax=3;m=2;}]
			\fill[pattern=north east lines] (xmin,ymin) rectangle (xmax,ymax);
			\fill[white] (xmin,ymin) rectangle (m,ymax);
			\draw[thick,dashed] (m,ymin) -- (m,ymax);
			\draw[->] (xmin,0) -- (xmax,0)node[above]{$x$};
			\foreach \x in {-1,1,2,3,4,5}
			\draw[shift={(\x,0)}] (0pt,2pt) -- (0pt,-2pt) node[below] {\footnotesize $\x$};
			\draw[->,color=black] (0,ymin) -- (0,ymax)node[left]{$y$};
			\foreach \y in {-1,1,2}
			\draw[shift={(0,\y)}] (2pt,0pt) -- (-2pt,0pt) node[left] {\footnotesize $\y$};
			\node[above left] at (0,0){$O$};
		\end{tikzpicture}}
	\loigiai{
		Đường thẳng bờ $d\colon x=2$ vẽ nét đứt nên không kể bờ.\\
		Lấy điểm $O(0;0)$ thuộc phần không bị gạch, ta có $0<2$ nên phần không bị gạch là miền nghiệm của bất phương trình $x<2$.
	}
\end{ex}

%%%==============Câu_5 — x = -2, tô bên trái, bờ NÉT ĐỨT=============%%%
\begin{ex}%[VNPmath - tự biên soạn]%[GV. Vũ Ngọc Phát]%[0D2N1-2]
	\immini[thm]{Phần không bị gạch (không kể đường thẳng $d$) trong hình vẽ là miền nghiệm của bất phương trình nào sau đây?
		\choice
		{$x<-3$}
		{$x>-3$}
		{$x<-2$}
		{\True $x>-2$}}
	{\begin{tikzpicture}[scale=.7,thick,>=stealth,declare function={xmin=-5;xmax=1;ymin=-2;ymax=3;m=-2;}]
			\fill[pattern=north east lines] (xmin,ymin) rectangle (xmax,ymax);
			\fill[white] (m,ymin) rectangle (xmax,ymax);
			\draw[thick,dashed] (m,ymin) -- (m,ymax);
			\draw[->] (xmin,0) -- (xmax,0)node[above]{$x$};
			\foreach \x in {-4,-3,-2,-1,1}
			\draw[shift={(\x,0)}] (0pt,2pt) -- (0pt,-2pt) node[below] {\footnotesize $\x$};
			\draw[->,color=black] (0,ymin) -- (0,ymax)node[left]{$y$};
			\foreach \y in {-1,1,2}
			\draw[shift={(0,\y)}] (2pt,0pt) -- (-2pt,0pt) node[left] {\footnotesize $\y$};
			\node[above left] at (0,0){$O$};
		\end{tikzpicture}}
	\loigiai{
		Đường thẳng bờ $d\colon x=-2$ vẽ nét đứt nên không kể bờ.\\
		Lấy điểm $O(0;0)$ thuộc phần không bị gạch, ta có $0>-2$ nên phần không bị gạch là miền nghiệm của bất phương trình $x>-2$.
	}
\end{ex}

%%%%%%===================================================================%%%%%%
%%%%%%============================= DẠNG 2 ===============================%%%%%%
%%%%%% Bờ là đường thẳng song song với trục hoành: y = m
%%%%%%===================================================================%%%%%%
\subsubsection{Bờ là đường thẳng $y=m$}

%%%==============Câu_6 — y = 0 (trục hoành), tô phía dưới, bờ NÉT LIỀN=============%%%
\begin{ex}%[VNPmath - tự biên soạn]%[GV. Vũ Ngọc Phát]%[0D2N1-2]
	\immini[thm]{Phần không bị gạch trong hình vẽ là miền nghiệm của bất phương trình nào sau đây?
		\choice
		{$y\leq 0$}
		{$x\leq 0$}
		{$x\geq 0$}
		{\True $y\geq 0$}}
	{\begin{tikzpicture}[scale=.7,thick,>=stealth,declare function={xmin=-2;xmax=4;ymin=-2.5;ymax=2.5;m=0;}]
			\fill[pattern=north east lines] (xmin,ymin) rectangle (xmax,ymax);
			\fill[white] (xmin,m) rectangle (xmax,ymax);
			\draw[->] (xmin,0) -- (xmax,0)node[above]{$x$};
			\foreach \x in {-1,1,2,3,4}
			\draw[shift={(\x,0)}] (0pt,2pt) -- (0pt,-2pt) node[below] {\footnotesize $\x$};
			\draw[->,color=black] (0,ymin) -- (0,ymax)node[left]{$y$};
			\foreach \y in {-2,-1,1,2}
			\draw[shift={(0,\y)}] (2pt,0pt) -- (-2pt,0pt) node[left] {\footnotesize $\y$};
			\node[above left] at (0,0){$O$};
		\end{tikzpicture}}
	\loigiai{
		Đường thẳng bờ là trục hoành $d\colon y=0$ (vẽ nét liền nên kể cả bờ).\\
		Lấy điểm $M(0;1)$ thuộc phần không bị gạch, ta có $1\geq 0$ nên phần không bị gạch là miền nghiệm của bất phương trình $y\geq 0$.
	}
\end{ex}

%%%==============Câu_7 — y = 1, tô phía dưới, bờ NÉT ĐỨT=============%%%
\begin{ex}%[VNPmath - tự biên soạn]%[GV. Vũ Ngọc Phát]%[0D2N1-2]
	\immini[thm]{Phần không bị gạch (không kể đường thẳng $d$) trong hình vẽ là miền nghiệm của bất phương trình nào sau đây?
		\choice
		{\True $y>1$}
		{$y<1$}
		{$y>2$}
		{$y<2$}}
	{\begin{tikzpicture}[scale=.7,thick,>=stealth,declare function={xmin=-2;xmax=4;ymin=-1.5;ymax=3.5;m=1;}]
			\fill[pattern=north east lines] (xmin,ymin) rectangle (xmax,ymax);
			\fill[white] (xmin,m) rectangle (xmax,ymax);
			\draw[thick,dashed] (xmin,m) -- (xmax,m);
			\draw[->] (xmin,0) -- (xmax,0)node[above]{$x$};
			\foreach \x in {-1,1,2,3,4}
			\draw[shift={(\x,0)}] (0pt,2pt) -- (0pt,-2pt) node[below] {\footnotesize $\x$};
			\draw[->,color=black] (0,ymin) -- (0,ymax)node[left]{$y$};
			\foreach \y in {-1,1,2,3}
			\draw[shift={(0,\y)}] (2pt,0pt) -- (-2pt,0pt) node[left] {\footnotesize $\y$};
			\node[above left] at (0,0){$O$};
		\end{tikzpicture}}
	\loigiai{
		Đường thẳng bờ $d\colon y=1$ vẽ nét đứt nên không kể bờ.\\
		Lấy điểm $M(0;2)$ thuộc phần không bị gạch, ta có $2>1$ nên phần không bị gạch là miền nghiệm của bất phương trình $y>1$.
	}
\end{ex}

%%%==============Câu_8 — y = 2, tô phía trên, bờ NÉT LIỀN=============%%%
\begin{ex}%[VNPmath - tự biên soạn]%[GV. Vũ Ngọc Phát]%[0D2N1-2]
	\immini[thm]{Phần không bị gạch trong hình vẽ là miền nghiệm của bất phương trình nào sau đây?
		\choice
		{$y\geq 2$}
		{$y\geq -2$}
		{\True $y\leq 2$}
		{$y\leq -2$}}
	{\begin{tikzpicture}[scale=.7,thick,>=stealth,declare function={xmin=-2;xmax=4;ymin=-0.5;ymax=4.5;m=2;}]
			\fill[pattern=north east lines] (xmin,ymin) rectangle (xmax,ymax);
			\fill[white] (xmin,ymin) rectangle (xmax,m);
			\draw[thick] (xmin,m) -- (xmax,m);
			\draw[->] (xmin,0) -- (xmax,0)node[above]{$x$};
			\foreach \x in {-1,1,2,3,4}
			\draw[shift={(\x,0)}] (0pt,2pt) -- (0pt,-2pt) node[below] {\footnotesize $\x$};
			\draw[->,color=black] (0,ymin) -- (0,ymax)node[left]{$y$};
			\foreach \y in {1,2,3,4}
			\draw[shift={(0,\y)}] (2pt,0pt) -- (-2pt,0pt) node[left] {\footnotesize $\y$};
			\node[above left] at (0,0){$O$};
		\end{tikzpicture}}
	\loigiai{
		Đường thẳng bờ $d\colon y=2$ (vẽ nét liền nên kể cả bờ).\\
		Lấy điểm $O(0;0)$ thuộc phần không bị gạch, ta có $0\leq 2$ nên phần không bị gạch là miền nghiệm của bất phương trình $y\leq 2$.
	}
\end{ex}

%%%==============Câu_9 — y = -1, tô phía trên, bờ NÉT ĐỨT=============%%%
\begin{ex}%[VNPmath - tự biên soạn]%[GV. Vũ Ngọc Phát]%[0D2N1-2]
	\immini[thm]{Phần không bị gạch (không kể đường thẳng $d$) trong hình vẽ là miền nghiệm của bất phương trình nào sau đây?
		\choice
		{$y>-1$}
		{\True $y<-1$}
		{$y>-2$}
		{$y<-2$}}
	{\begin{tikzpicture}[scale=.7,thick,>=stealth,declare function={xmin=-2;xmax=4;ymin=-3.5;ymax=1.5;m=-1;}]
			\fill[pattern=north east lines] (xmin,ymin) rectangle (xmax,ymax);
			\fill[white] (xmin,ymin) rectangle (xmax,m);
			\draw[thick,dashed] (xmin,m) -- (xmax,m);
			\draw[->] (xmin,0) -- (xmax,0)node[above]{$x$};
			\foreach \x in {-1,1,2,3,4}
			\draw[shift={(\x,0)}] (0pt,2pt) -- (0pt,-2pt) node[below] {\footnotesize $\x$};
			\draw[->,color=black] (0,ymin) -- (0,ymax)node[left]{$y$};
			\foreach \y in {-3,-2,-1,1}
			\draw[shift={(0,\y)}] (2pt,0pt) -- (-2pt,0pt) node[left] {\footnotesize $\y$};
			\node[above left] at (0,0){$O$};
		\end{tikzpicture}}
	\loigiai{
		Đường thẳng bờ $d\colon y=-1$ vẽ nét đứt nên không kể bờ.\\
		Lấy điểm $M(0;-2)$ thuộc phần không bị gạch, ta có $-2<-1$ nên phần không bị gạch là miền nghiệm của bất phương trình $y<-1$.
	}
\end{ex}

%%%%%%===================================================================%%%%%%
%%%%%%============================= DẠNG 3 ===============================%%%%%%
%%%%%% Bờ là đường thẳng ax+by = c với ab < 0  (hệ số góc dương)
%%%%%%===================================================================%%%%%%
\subsubsection{Bờ là đường thẳng $ax+by=c$ với $ab<0$}

%%%==============Câu_10 — y = x (x - y = 0), tô nửa trên - trái, bờ NÉT LIỀN=============%%%
\begin{ex}%[VNPmath - tự biên soạn]%[GV. Vũ Ngọc Phát]%[0D2H1-2]
	\immini[thm]{Phần không bị gạch trong hình vẽ là miền nghiệm của bất phương trình nào sau đây?
		\choice
		{$y\geq x$}
		{$y\geq -x$}
		{$y\leq -x$}
		{\True $y\leq x$}}
	{\begin{tikzpicture}[scale=.7,thick,>=stealth,declare function={xmin=-2.5;xmax=3.5;ymin=-2;ymax=3;f(\x)=\x;}]
			\begin{scope}
				\clip (xmin,ymin) rectangle (xmax,ymax);
				\fill[pattern=north east lines] (xmin,ymin) rectangle (xmax,ymax);
				\fill[white] (xmin,{f(xmin)}) -- (xmax,{f(xmax)}) -- (xmax,ymin) -- (xmin,ymin) -- cycle;
				\draw[thick] (xmin,{f(xmin)}) -- (xmax,{f(xmax)});
			\end{scope}
			\draw[->] (xmin,0) -- (xmax,0)node[above]{$x$};
			\foreach \x in {-2,-1,1,2,3}
			\draw[shift={(\x,0)}] (0pt,2pt) -- (0pt,-2pt) node[below] {\footnotesize $\x$};
			\draw[->,color=black] (0,ymin) -- (0,ymax)node[left]{$y$};
			\foreach \y in {-1,1,2}
			\draw[shift={(0,\y)}] (2pt,0pt) -- (-2pt,0pt) node[left] {\footnotesize $\y$};
			\node[above left] at (0,0){$O$};
		\end{tikzpicture}}
	\loigiai{
		Đường thẳng bờ đi qua gốc toạ độ $O$ và điểm $(1;1)$ nên có phương trình $y=x$ (vẽ nét liền nên kể cả bờ).\\
		Lấy điểm $M(1;0)$ thuộc phần không bị gạch, ta có $0\leq 1$ nên phần không bị gạch là miền nghiệm của bất phương trình $y\leq x$.
	}
\end{ex}

%%%==============Câu_11 — y = 2x (2x - y = 0), tô nửa trên - trái, bờ NÉT LIỀN=============%%%
\begin{ex}%[VNPmath - tự biên soạn]%[GV. Vũ Ngọc Phát]%[0D2H1-2]
	\immini[thm]{Phần không bị gạch trong hình vẽ là miền nghiệm của bất phương trình nào sau đây?
		\choice
		{$y\geq 2x$}
		{$2y\geq x$}
		{$2y\leq x$}
		{\True $y\leq 2x$}}
	{\begin{tikzpicture}[scale=.7,thick,>=stealth,declare function={xmin=-2.5;xmax=3.5;ymin=-1.5;ymax=3.5;f(\x)=2*(\x);}]
			\begin{scope}
				\clip (xmin,ymin) rectangle (xmax,ymax);
				\fill[pattern=north east lines] (xmin,ymin) rectangle (xmax,ymax);
				\fill[white] (xmin,{f(xmin)}) -- (xmax,{f(xmax)}) -- (xmax,ymin) -- (xmin,ymin) -- cycle;
				\draw[thick] (xmin,{f(xmin)}) -- (xmax,{f(xmax)});
			\end{scope}
			\draw[->] (xmin,0) -- (xmax,0)node[above]{$x$};
			\foreach \x in {-2,-1,1,2,3}
			\draw[shift={(\x,0)}] (0pt,2pt) -- (0pt,-2pt) node[below] {\footnotesize $\x$};
			\draw[->,color=black] (0,ymin) -- (0,ymax)node[left]{$y$};
			\foreach \y in {-1,1,2,3}
			\draw[shift={(0,\y)}] (2pt,0pt) -- (-2pt,0pt) node[left] {\footnotesize $\y$};
			\node[above left] at (0,0){$O$};
		\end{tikzpicture}}
	\loigiai{
		Đường thẳng bờ đi qua gốc toạ độ $O$ và điểm $(1;2)$ nên có phương trình $y=2x$ (vẽ nét liền nên kể cả bờ).\\
		Lấy điểm $M(1;0)$ thuộc phần không bị gạch, ta có $0\leq 2$ nên phần không bị gạch là miền nghiệm của bất phương trình $y\leq 2x$.
	}
\end{ex}

%%%==============Câu_12 — y = x/2 + 1 (x - 2y + 2 = 0), tô nửa dưới - phải, bờ NÉT ĐỨT=============%%%
\begin{ex}%[VNPmath - tự biên soạn]%[GV. Vũ Ngọc Phát]%[0D2H1-2]
	\immini[thm]{Phần không bị gạch (không kể đường thẳng $d$) trong hình vẽ là miền nghiệm của bất phương trình nào sau đây?
		\choice
		{\True $x-2y+2<0$}
		{$x-2y+2>0$}
		{$2x-y+2<0$}
		{$2x-y+2>0$}}
	{\begin{tikzpicture}[scale=.7,thick,>=stealth,declare function={xmin=-3;xmax=3;ymin=-1.5;ymax=3.5;f(\x)=1+(\x)/2;}]
			\begin{scope}
				\clip (xmin,ymin) rectangle (xmax,ymax);
				\fill[pattern=north east lines] (xmin,ymin) rectangle (xmax,ymax);
				\fill[white] (xmin,{f(xmin)}) -- (xmax,{f(xmax)}) -- (xmax,ymax) -- (xmin,ymax) -- cycle;
				\draw[thick,dashed] (xmin,{f(xmin)}) -- (xmax,{f(xmax)});
			\end{scope}
			\draw[->] (xmin,0) -- (xmax,0)node[above]{$x$};
			\foreach \x in {-2,-1,1,2,3}
			\draw[shift={(\x,0)}] (0pt,2pt) -- (0pt,-2pt) node[below] {\footnotesize $\x$};
			\draw[->,color=black] (0,ymin) -- (0,ymax)node[left]{$y$};
			\foreach \y in {-1,1,2,3}
			\draw[shift={(0,\y)}] (2pt,0pt) -- (-2pt,0pt) node[left] {\footnotesize $\y$};
			\node[above left] at (0,0){$O$};
		\end{tikzpicture}}
	\loigiai{
		Đường thẳng bờ $d$ đi qua hai điểm $(0;1)$ và $(2;2)$ nên có phương trình $y=\dfrac{x}{2}+1$, hay $x-2y+2=0$.\\
		$d$ vẽ nét đứt nên không kể bờ.\\
		Lấy điểm $M(0;2)$ thuộc phần không bị gạch, ta có $0-2\cdot 2+2=-2<0$ nên phần không bị gạch là miền nghiệm của bất phương trình $x-2y+2<0$.
	}
\end{ex}

%%%==============Câu_13 — y = x + 2 (x - y + 2 = 0), tô nửa trên - trái, bờ NÉT ĐỨT=============%%%
\begin{ex}%[VNPmath - tự biên soạn]%[GV. Vũ Ngọc Phát]%[0D2H1-2]
	\immini[thm]{Phần không bị gạch (không kể đường thẳng $d$) trong hình vẽ là miền nghiệm của bất phương trình nào sau đây?
		\choice
		{$x-y+2<0$}
		{\True $x-y+2>0$}
		{$x+y+2<0$}
		{$x+y+2>0$}}
	{\begin{tikzpicture}[scale=.7,thick,>=stealth,declare function={xmin=-4;xmax=2;ymin=-1.5;ymax=3.5;f(\x)=(\x)+2;}]
			\begin{scope}
				\clip (xmin,ymin) rectangle (xmax,ymax);
				\fill[pattern=north east lines] (xmin,ymin) rectangle (xmax,ymax);
				\fill[white] (xmin,{f(xmin)}) -- (xmax,{f(xmax)}) -- (xmax,ymin) -- (xmin,ymin) -- cycle;
				\draw[thick,dashed] (xmin,{f(xmin)}) -- (xmax,{f(xmax)});
			\end{scope}
			\draw[->] (xmin,0) -- (xmax,0)node[above]{$x$};
			\foreach \x in {-3,-2,-1,1,2}
			\draw[shift={(\x,0)}] (0pt,2pt) -- (0pt,-2pt) node[below] {\footnotesize $\x$};
			\draw[->,color=black] (0,ymin) -- (0,ymax)node[left]{$y$};
			\foreach \y in {-1,1,2,3}
			\draw[shift={(0,\y)}] (2pt,0pt) -- (-2pt,0pt) node[left] {\footnotesize $\y$};
			\node[above left] at (0,0){$O$};
		\end{tikzpicture}}
	\loigiai{
		Đường thẳng bờ $d$ đi qua hai điểm $(-2;0)$ và $(0;2)$ nên có phương trình $y=x+2$, hay $x-y+2=0$.\\
		$d$ vẽ nét đứt nên không kể bờ.\\
		Lấy điểm $M(0;1)$ thuộc phần không bị gạch, ta có $0-1+2=1>0$ nên phần không bị gạch là miền nghiệm của bất phương trình $x-y+2>0$.
	}
\end{ex}

%%%%%%===================================================================%%%%%%
%%%%%%============================= DẠNG 4 ===============================%%%%%%
%%%%%% Bờ là đường thẳng ax+by = c với ab > 0  (hệ số góc âm)
%%%%%%===================================================================%%%%%%
\subsubsection{Bờ là đường thẳng $ax+by=c$ với $ab>0$}

%%%==============Câu_14 — x + y = 1, tô nửa dưới - trái, bờ NÉT ĐỨT=============%%%
\begin{ex}%[VNPmath - tự biên soạn]%[GV. Vũ Ngọc Phát]%[0D2H1-2]
	\immini[thm]{Phần không bị gạch (không kể đường thẳng $d$) trong hình vẽ là miền nghiệm của bất phương trình nào sau đây?
		\choice
		{$x+y<1$}
		{$x-y>1$}
		{\True $x+y>1$}
		{$x-y<1$}}
	{\begin{tikzpicture}[scale=.7,thick,>=stealth,declare function={xmin=-2.5;xmax=3.5;ymin=-2;ymax=3;f(\x)=1-(\x);}]
			\begin{scope}
				\clip (xmin,ymin) rectangle (xmax,ymax);
				\fill[pattern=north east lines] (xmin,ymin) rectangle (xmax,ymax);
				\fill[white] (xmin,{f(xmin)}) -- (xmax,{f(xmax)}) -- (xmax,ymax) -- (xmin,ymax) -- cycle;
				\draw[thick,dashed] (xmin,{f(xmin)}) -- (xmax,{f(xmax)});
			\end{scope}
			\draw[->] (xmin,0) -- (xmax,0)node[above]{$x$};
			\foreach \x in {-2,-1,1,2,3}
			\draw[shift={(\x,0)}] (0pt,2pt) -- (0pt,-2pt) node[below] {\footnotesize $\x$};
			\draw[->,color=black] (0,ymin) -- (0,ymax)node[left]{$y$};
			\foreach \y in {-1,1,2}
			\draw[shift={(0,\y)}] (2pt,0pt) -- (-2pt,0pt) node[left] {\footnotesize $\y$};
			\node[above left] at (0,0){$O$};
		\end{tikzpicture}}
	\loigiai{
		Đường thẳng bờ $d$ đi qua hai điểm $(1;0)$ và $(0;1)$ nên có phương trình $x+y=1$.\\
		$d$ vẽ nét đứt nên không kể bờ.\\
		Lấy điểm $M(2;0)$ thuộc phần không bị gạch, ta có $2+0=2>1$ nên phần không bị gạch là miền nghiệm của bất phương trình $x+y>1$.
	}
\end{ex}

%%%==============Câu_15 — x + y = 3, tô nửa dưới - trái, bờ NÉT LIỀN=============%%%
\begin{ex}%[VNPmath - tự biên soạn]%[GV. Vũ Ngọc Phát]%[0D2H1-2]
	\immini[thm]{Phần không bị gạch trong hình vẽ là miền nghiệm của bất phương trình nào sau đây?
		\choice
		{$x+y<3$}
		{\True $x+y\geq 3$}
		{$x-y\geq 3$}
		{$x+y\leq 3$}}
	{\begin{tikzpicture}[scale=.7,thick,>=stealth,declare function={xmin=-1.5;xmax=4.5;ymin=-1;ymax=4;f(\x)=3-(\x);}]
			\begin{scope}
				\clip (xmin,ymin) rectangle (xmax,ymax);
				\fill[pattern=north east lines] (xmin,ymin) rectangle (xmax,ymax);
				\fill[white] (xmin,{f(xmin)}) -- (xmax,{f(xmax)}) -- (xmax,ymax) -- (xmin,ymax) -- cycle;
				\draw[thick] (xmin,{f(xmin)}) -- (xmax,{f(xmax)});
			\end{scope}
			\draw[->] (xmin,0) -- (xmax,0)node[above]{$x$};
			\foreach \x in {-1,1,2,3,4}
			\draw[shift={(\x,0)}] (0pt,2pt) -- (0pt,-2pt) node[below] {\footnotesize $\x$};
			\draw[->,color=black] (0,ymin) -- (0,ymax)node[left]{$y$};
			\foreach \y in {-1,1,2,3}
			\draw[shift={(0,\y)}] (2pt,0pt) -- (-2pt,0pt) node[left] {\footnotesize $\y$};
			\node[above left] at (0,0){$O$};
		\end{tikzpicture}}
	\loigiai{
		Đường thẳng bờ đi qua hai điểm $(3;0)$ và $(0;3)$ nên có phương trình $x+y=3$ (vẽ nét liền nên kể cả bờ).\\
		Lấy điểm $M(4;0)$ thuộc phần không bị gạch, ta có $4+0=4\geq 3$ nên phần không bị gạch là miền nghiệm của bất phương trình $x+y\geq 3$.
	}
\end{ex}

%%%==============Câu_16 — x + 2y = 2, tô nửa trên - phải, bờ NÉT ĐỨT=============%%%
\begin{ex}%[VNPmath - tự biên soạn]%[GV. Vũ Ngọc Phát]%[0D2H1-2]
	\immini[thm]{Phần không bị gạch (không kể đường thẳng $d$) trong hình vẽ là miền nghiệm của bất phương trình nào sau đây?
		\choice
		{$x+2y>2$}
		{$2x+y>2$}
		{\True $x+2y<2$}
		{$2x+y<2$}}
	{\begin{tikzpicture}[scale=.7,thick,>=stealth,declare function={xmin=-2;xmax=4;ymin=-2;ymax=3;f(\x)=1-(\x)/2;}]
			\begin{scope}
				\clip (xmin,ymin) rectangle (xmax,ymax);
				\fill[pattern=north east lines] (xmin,ymin) rectangle (xmax,ymax);
				\fill[white] (xmin,{f(xmin)}) -- (xmax,{f(xmax)}) -- (xmax,ymin) -- (xmin,ymin) -- cycle;
				\draw[thick,dashed] (xmin,{f(xmin)}) -- (xmax,{f(xmax)});
			\end{scope}
			\draw[->] (xmin,0) -- (xmax,0)node[above]{$x$};
			\foreach \x in {-1,1,2,3,4}
			\draw[shift={(\x,0)}] (0pt,2pt) -- (0pt,-2pt) node[below] {\footnotesize $\x$};
			\draw[->,color=black] (0,ymin) -- (0,ymax)node[left]{$y$};
			\foreach \y in {-1,1,2}
			\draw[shift={(0,\y)}] (2pt,0pt) -- (-2pt,0pt) node[left] {\footnotesize $\y$};
			\node[above left] at (0,0){$O$};
		\end{tikzpicture}}
	\loigiai{
		Đường thẳng bờ $d$ đi qua hai điểm $(0;1)$ và $(2;0)$ nên có phương trình $y=1-\dfrac{x}{2}$, hay $x+2y=2$.\\
		$d$ vẽ nét đứt nên không kể bờ.\\
		Lấy điểm $O(0;0)$ thuộc phần không bị gạch, ta có $0+2\cdot 0=0<2$ nên phần không bị gạch là miền nghiệm của bất phương trình $x+2y<2$.
	}
\end{ex}


\Closesolutionfile{ans}

% 2) Nếu muốn GỘP vào TN-0D2-B1.tex: xoá 2 dòng \Opensolutionfile/\Closesolutionfile
%    ở file này rồi dán phần các câu hỏi vào TRƯỚC dòng \Closesolutionfile{ans} của file đó.
% 3) In bảng đáp án: \inputansbox{10}{ans/ans-0D2-B1-MienNghiem}
% 4) Lời giải ngắn đã điền sẵn trong \loigiai{}; xoá nội dung nếu muốn để trống như mẫu gốc.
%    QUY ƯỚC LỜI GIẢI: chỉ phân tích miền KHÔNG bị gạch — nêu phương trình đường bờ (kèm ghi
%    chú nét liền/nét đứt), rồi lấy điểm $M$ thuộc phần không bị gạch (ưu tiên điểm có ít nhất
%    một toạ độ bằng 0) thay vào bất phương trình để kết luận.
%    - Nếu lấy gốc toạ độ thì gọi là $O(0;0)$ (KHÔNG dùng $M(0;0)$).
% 5) Đang dùng \subsubsection làm tiêu đề dạng (đánh số 1., 2., 3., 4.; KHÔNG reset số câu,
%    nên bảng đáp án vẫn liên tục 1 → 16). Nếu muốn khung "DẠNG TOÁN" thì thay bằng
%    \begin{dang}{...} ... \end{dang} — lưu ý môi trường dang tự reset số câu về 1.
% 6) Quy ước nhãn câu: %[Nguồn]%[Người biên soạn]%[Mã ID].
%    Mã ID dùng ở đây: 0D2N1-2 (nhận biết — bờ song song trục) và 0D2H1-2 (thông hiểu — bờ xiên),
%    khớp với data/10/0D2/0-KNTT-C2B3-X.tex.
%
% GHI CHÚ: hai câu bờ x = m, y = m dạng nét đứt trước đây dùng m = 3 (khung phải lệch, khó đạt
% 50% mà vẫn thấy gốc O) đã đổi thành m = -2 (câu 5) và m = -1 (câu 9) để khung đối xứng qua d.
%%%%%%%%%%%%%%%%%%%%%%%%%%%%%%%%%%%%%%%%%%%%%%%%%%%%%%%%%%%%%%%%%%%%%%%%%%%%%%%%%%%%%

\Opensolutionfile{ans}[ans/ans-0D2-B1-MienNghiem]

%%%%%%===================================================================%%%%%%
%%%%%%============================= DẠNG 1 ===============================%%%%%%
%%%%%% Bờ là đường thẳng song song với trục tung: x = m
%%%%%%===================================================================%%%%%%
\subsubsection{Bờ là đường thẳng $x=m$}

%%%==============Câu_1 — x = 0 (trục tung), tô bên trái, bờ NÉT LIỀN=============%%%
\begin{ex}%[VNPmath - tự biên soạn]%[GV. Vũ Ngọc Phát]%[0D2N1-2]
	\immini[thm]{Phần không bị gạch trong hình vẽ là miền nghiệm của bất phương trình nào sau đây?
		\choice
		{$x\leq 0$}
		{\True $x\geq 0$}
		{$y\leq 0$}
		{$y\geq 0$}}
	{\begin{tikzpicture}[scale=.7,thick,>=stealth,declare function={xmin=-3;xmax=3;ymin=-2;ymax=3;m=0;}]
			\fill[pattern=north east lines] (xmin,ymin) rectangle (xmax,ymax);
			\fill[white] (m,ymin) rectangle (xmax,ymax);
			\draw[->] (xmin,0) -- (xmax,0)node[above]{$x$};
			\foreach \x in {-2,-1,1,2,3}
			\draw[shift={(\x,0)}] (0pt,2pt) -- (0pt,-2pt) node[below] {\footnotesize $\x$};
			\draw[->,color=black] (0,ymin) -- (0,ymax)node[left]{$y$};
			\foreach \y in {-1,1,2}
			\draw[shift={(0,\y)}] (2pt,0pt) -- (-2pt,0pt) node[left] {\footnotesize $\y$};
			\node[above left] at (0,0){$O$};
		\end{tikzpicture}}
	\loigiai{
		Đường thẳng bờ là trục tung $d\colon x=0$ (vẽ nét liền nên kể cả bờ).\\
		Lấy điểm $M(1;0)$ thuộc phần không bị gạch, ta có $1\geq 0$ nên phần không bị gạch là miền nghiệm của bất phương trình $x\geq 0$.
	}
\end{ex}

%%%==============Câu_2 — x = 1, tô bên trái, bờ NÉT LIỀN=============%%%
\begin{ex}%[VNPmath - tự biên soạn]%[GV. Vũ Ngọc Phát]%[0D2N1-2]
	\immini[thm]{Phần không bị gạch trong hình vẽ là miền nghiệm của bất phương trình nào sau đây?
		\choice
		{\True $x\geq 1$}
		{$x\leq 1$}
		{$x\geq -1$}
		{$x\leq -1$}}
	{\begin{tikzpicture}[scale=.7,thick,>=stealth,declare function={xmin=-2;xmax=4;ymin=-2;ymax=3;m=1;}]
			\fill[pattern=north east lines] (xmin,ymin) rectangle (xmax,ymax);
			\fill[white] (m,ymin) rectangle (xmax,ymax);
			\draw[thick] (m,ymin) -- (m,ymax);
			\draw[->] (xmin,0) -- (xmax,0)node[above]{$x$};
			\foreach \x in {-1,1,2,3,4}
			\draw[shift={(\x,0)}] (0pt,2pt) -- (0pt,-2pt) node[below] {\footnotesize $\x$};
			\draw[->,color=black] (0,ymin) -- (0,ymax)node[left]{$y$};
			\foreach \y in {-1,1,2}
			\draw[shift={(0,\y)}] (2pt,0pt) -- (-2pt,0pt) node[left] {\footnotesize $\y$};
			\node[above left] at (0,0){$O$};
		\end{tikzpicture}}
	\loigiai{
		Đường thẳng bờ $d\colon x=1$ (vẽ nét liền nên kể cả bờ).\\
		Lấy điểm $M(2;0)$ thuộc phần không bị gạch, ta có $2\geq 1$ nên phần không bị gạch là miền nghiệm của bất phương trình $x\geq 1$.
	}
\end{ex}

%%%==============Câu_3 — x = -1, tô bên phải, bờ NÉT LIỀN=============%%%
\begin{ex}%[VNPmath - tự biên soạn]%[GV. Vũ Ngọc Phát]%[0D2N1-2]
	\immini[thm]{Phần không bị gạch trong hình vẽ là miền nghiệm của bất phương trình nào sau đây?
		\choice
		{$x\geq -1$}
		{$x\geq 1$}
		{\True $x\leq -1$}
		{$x\leq 1$}}
	{\begin{tikzpicture}[scale=.7,thick,>=stealth,declare function={xmin=-4;xmax=2;ymin=-2;ymax=3;m=-1;}]
			\fill[pattern=north east lines] (xmin,ymin) rectangle (xmax,ymax);
			\fill[white] (xmin,ymin) rectangle (m,ymax);
			\draw[thick] (m,ymin) -- (m,ymax);
			\draw[->] (xmin,0) -- (xmax,0)node[above]{$x$};
			\foreach \x in {-3,-2,-1,1,2}
			\draw[shift={(\x,0)}] (0pt,2pt) -- (0pt,-2pt) node[below] {\footnotesize $\x$};
			\draw[->,color=black] (0,ymin) -- (0,ymax)node[left]{$y$};
			\foreach \y in {-1,1,2}
			\draw[shift={(0,\y)}] (2pt,0pt) -- (-2pt,0pt) node[left] {\footnotesize $\y$};
			\node[above left] at (0,0){$O$};
		\end{tikzpicture}}
	\loigiai{
		Đường thẳng bờ $d\colon x=-1$ (vẽ nét liền nên kể cả bờ).\\
		Lấy điểm $M(-2;0)$ thuộc phần không bị gạch, ta có $-2\leq -1$ nên phần không bị gạch là miền nghiệm của bất phương trình $x\leq -1$.
	}
\end{ex}

%%%==============Câu_4 — x = 2, tô bên phải, bờ NÉT ĐỨT=============%%%
\begin{ex}%[VNPmath - tự biên soạn]%[GV. Vũ Ngọc Phát]%[0D2N1-2]
	\immini[thm]{Phần không bị gạch (không kể đường thẳng $d$) trong hình vẽ là miền nghiệm của bất phương trình nào sau đây?
		\choice
		{$x>2$}
		{$x<3$}
		{$x>3$}
		{\True $x<2$}}
	{\begin{tikzpicture}[scale=.7,thick,>=stealth,declare function={xmin=-1;xmax=5;ymin=-2;ymax=3;m=2;}]
			\fill[pattern=north east lines] (xmin,ymin) rectangle (xmax,ymax);
			\fill[white] (xmin,ymin) rectangle (m,ymax);
			\draw[thick,dashed] (m,ymin) -- (m,ymax);
			\draw[->] (xmin,0) -- (xmax,0)node[above]{$x$};
			\foreach \x in {-1,1,2,3,4,5}
			\draw[shift={(\x,0)}] (0pt,2pt) -- (0pt,-2pt) node[below] {\footnotesize $\x$};
			\draw[->,color=black] (0,ymin) -- (0,ymax)node[left]{$y$};
			\foreach \y in {-1,1,2}
			\draw[shift={(0,\y)}] (2pt,0pt) -- (-2pt,0pt) node[left] {\footnotesize $\y$};
			\node[above left] at (0,0){$O$};
		\end{tikzpicture}}
	\loigiai{
		Đường thẳng bờ $d\colon x=2$ vẽ nét đứt nên không kể bờ.\\
		Lấy điểm $O(0;0)$ thuộc phần không bị gạch, ta có $0<2$ nên phần không bị gạch là miền nghiệm của bất phương trình $x<2$.
	}
\end{ex}

%%%==============Câu_5 — x = -2, tô bên trái, bờ NÉT ĐỨT=============%%%
\begin{ex}%[VNPmath - tự biên soạn]%[GV. Vũ Ngọc Phát]%[0D2N1-2]
	\immini[thm]{Phần không bị gạch (không kể đường thẳng $d$) trong hình vẽ là miền nghiệm của bất phương trình nào sau đây?
		\choice
		{$x<-3$}
		{$x>-3$}
		{$x<-2$}
		{\True $x>-2$}}
	{\begin{tikzpicture}[scale=.7,thick,>=stealth,declare function={xmin=-5;xmax=1;ymin=-2;ymax=3;m=-2;}]
			\fill[pattern=north east lines] (xmin,ymin) rectangle (xmax,ymax);
			\fill[white] (m,ymin) rectangle (xmax,ymax);
			\draw[thick,dashed] (m,ymin) -- (m,ymax);
			\draw[->] (xmin,0) -- (xmax,0)node[above]{$x$};
			\foreach \x in {-4,-3,-2,-1,1}
			\draw[shift={(\x,0)}] (0pt,2pt) -- (0pt,-2pt) node[below] {\footnotesize $\x$};
			\draw[->,color=black] (0,ymin) -- (0,ymax)node[left]{$y$};
			\foreach \y in {-1,1,2}
			\draw[shift={(0,\y)}] (2pt,0pt) -- (-2pt,0pt) node[left] {\footnotesize $\y$};
			\node[above left] at (0,0){$O$};
		\end{tikzpicture}}
	\loigiai{
		Đường thẳng bờ $d\colon x=-2$ vẽ nét đứt nên không kể bờ.\\
		Lấy điểm $O(0;0)$ thuộc phần không bị gạch, ta có $0>-2$ nên phần không bị gạch là miền nghiệm của bất phương trình $x>-2$.
	}
\end{ex}

%%%%%%===================================================================%%%%%%
%%%%%%============================= DẠNG 2 ===============================%%%%%%
%%%%%% Bờ là đường thẳng song song với trục hoành: y = m
%%%%%%===================================================================%%%%%%
\subsubsection{Bờ là đường thẳng $y=m$}

%%%==============Câu_6 — y = 0 (trục hoành), tô phía dưới, bờ NÉT LIỀN=============%%%
\begin{ex}%[VNPmath - tự biên soạn]%[GV. Vũ Ngọc Phát]%[0D2N1-2]
	\immini[thm]{Phần không bị gạch trong hình vẽ là miền nghiệm của bất phương trình nào sau đây?
		\choice
		{$y\leq 0$}
		{$x\leq 0$}
		{$x\geq 0$}
		{\True $y\geq 0$}}
	{\begin{tikzpicture}[scale=.7,thick,>=stealth,declare function={xmin=-2;xmax=4;ymin=-2.5;ymax=2.5;m=0;}]
			\fill[pattern=north east lines] (xmin,ymin) rectangle (xmax,ymax);
			\fill[white] (xmin,m) rectangle (xmax,ymax);
			\draw[->] (xmin,0) -- (xmax,0)node[above]{$x$};
			\foreach \x in {-1,1,2,3,4}
			\draw[shift={(\x,0)}] (0pt,2pt) -- (0pt,-2pt) node[below] {\footnotesize $\x$};
			\draw[->,color=black] (0,ymin) -- (0,ymax)node[left]{$y$};
			\foreach \y in {-2,-1,1,2}
			\draw[shift={(0,\y)}] (2pt,0pt) -- (-2pt,0pt) node[left] {\footnotesize $\y$};
			\node[above left] at (0,0){$O$};
		\end{tikzpicture}}
	\loigiai{
		Đường thẳng bờ là trục hoành $d\colon y=0$ (vẽ nét liền nên kể cả bờ).\\
		Lấy điểm $M(0;1)$ thuộc phần không bị gạch, ta có $1\geq 0$ nên phần không bị gạch là miền nghiệm của bất phương trình $y\geq 0$.
	}
\end{ex}

%%%==============Câu_7 — y = 1, tô phía dưới, bờ NÉT ĐỨT=============%%%
\begin{ex}%[VNPmath - tự biên soạn]%[GV. Vũ Ngọc Phát]%[0D2N1-2]
	\immini[thm]{Phần không bị gạch (không kể đường thẳng $d$) trong hình vẽ là miền nghiệm của bất phương trình nào sau đây?
		\choice
		{\True $y>1$}
		{$y<1$}
		{$y>2$}
		{$y<2$}}
	{\begin{tikzpicture}[scale=.7,thick,>=stealth,declare function={xmin=-2;xmax=4;ymin=-1.5;ymax=3.5;m=1;}]
			\fill[pattern=north east lines] (xmin,ymin) rectangle (xmax,ymax);
			\fill[white] (xmin,m) rectangle (xmax,ymax);
			\draw[thick,dashed] (xmin,m) -- (xmax,m);
			\draw[->] (xmin,0) -- (xmax,0)node[above]{$x$};
			\foreach \x in {-1,1,2,3,4}
			\draw[shift={(\x,0)}] (0pt,2pt) -- (0pt,-2pt) node[below] {\footnotesize $\x$};
			\draw[->,color=black] (0,ymin) -- (0,ymax)node[left]{$y$};
			\foreach \y in {-1,1,2,3}
			\draw[shift={(0,\y)}] (2pt,0pt) -- (-2pt,0pt) node[left] {\footnotesize $\y$};
			\node[above left] at (0,0){$O$};
		\end{tikzpicture}}
	\loigiai{
		Đường thẳng bờ $d\colon y=1$ vẽ nét đứt nên không kể bờ.\\
		Lấy điểm $M(0;2)$ thuộc phần không bị gạch, ta có $2>1$ nên phần không bị gạch là miền nghiệm của bất phương trình $y>1$.
	}
\end{ex}

%%%==============Câu_8 — y = 2, tô phía trên, bờ NÉT LIỀN=============%%%
\begin{ex}%[VNPmath - tự biên soạn]%[GV. Vũ Ngọc Phát]%[0D2N1-2]
	\immini[thm]{Phần không bị gạch trong hình vẽ là miền nghiệm của bất phương trình nào sau đây?
		\choice
		{$y\geq 2$}
		{$y\geq -2$}
		{\True $y\leq 2$}
		{$y\leq -2$}}
	{\begin{tikzpicture}[scale=.7,thick,>=stealth,declare function={xmin=-2;xmax=4;ymin=-0.5;ymax=4.5;m=2;}]
			\fill[pattern=north east lines] (xmin,ymin) rectangle (xmax,ymax);
			\fill[white] (xmin,ymin) rectangle (xmax,m);
			\draw[thick] (xmin,m) -- (xmax,m);
			\draw[->] (xmin,0) -- (xmax,0)node[above]{$x$};
			\foreach \x in {-1,1,2,3,4}
			\draw[shift={(\x,0)}] (0pt,2pt) -- (0pt,-2pt) node[below] {\footnotesize $\x$};
			\draw[->,color=black] (0,ymin) -- (0,ymax)node[left]{$y$};
			\foreach \y in {1,2,3,4}
			\draw[shift={(0,\y)}] (2pt,0pt) -- (-2pt,0pt) node[left] {\footnotesize $\y$};
			\node[above left] at (0,0){$O$};
		\end{tikzpicture}}
	\loigiai{
		Đường thẳng bờ $d\colon y=2$ (vẽ nét liền nên kể cả bờ).\\
		Lấy điểm $O(0;0)$ thuộc phần không bị gạch, ta có $0\leq 2$ nên phần không bị gạch là miền nghiệm của bất phương trình $y\leq 2$.
	}
\end{ex}

%%%==============Câu_9 — y = -1, tô phía trên, bờ NÉT ĐỨT=============%%%
\begin{ex}%[VNPmath - tự biên soạn]%[GV. Vũ Ngọc Phát]%[0D2N1-2]
	\immini[thm]{Phần không bị gạch (không kể đường thẳng $d$) trong hình vẽ là miền nghiệm của bất phương trình nào sau đây?
		\choice
		{$y>-1$}
		{\True $y<-1$}
		{$y>-2$}
		{$y<-2$}}
	{\begin{tikzpicture}[scale=.7,thick,>=stealth,declare function={xmin=-2;xmax=4;ymin=-3.5;ymax=1.5;m=-1;}]
			\fill[pattern=north east lines] (xmin,ymin) rectangle (xmax,ymax);
			\fill[white] (xmin,ymin) rectangle (xmax,m);
			\draw[thick,dashed] (xmin,m) -- (xmax,m);
			\draw[->] (xmin,0) -- (xmax,0)node[above]{$x$};
			\foreach \x in {-1,1,2,3,4}
			\draw[shift={(\x,0)}] (0pt,2pt) -- (0pt,-2pt) node[below] {\footnotesize $\x$};
			\draw[->,color=black] (0,ymin) -- (0,ymax)node[left]{$y$};
			\foreach \y in {-3,-2,-1,1}
			\draw[shift={(0,\y)}] (2pt,0pt) -- (-2pt,0pt) node[left] {\footnotesize $\y$};
			\node[above left] at (0,0){$O$};
		\end{tikzpicture}}
	\loigiai{
		Đường thẳng bờ $d\colon y=-1$ vẽ nét đứt nên không kể bờ.\\
		Lấy điểm $M(0;-2)$ thuộc phần không bị gạch, ta có $-2<-1$ nên phần không bị gạch là miền nghiệm của bất phương trình $y<-1$.
	}
\end{ex}

%%%%%%===================================================================%%%%%%
%%%%%%============================= DẠNG 3 ===============================%%%%%%
%%%%%% Bờ là đường thẳng ax+by = c với ab < 0  (hệ số góc dương)
%%%%%%===================================================================%%%%%%
\subsubsection{Bờ là đường thẳng $ax+by=c$ với $ab<0$}

%%%==============Câu_10 — y = x (x - y = 0), tô nửa trên - trái, bờ NÉT LIỀN=============%%%
\begin{ex}%[VNPmath - tự biên soạn]%[GV. Vũ Ngọc Phát]%[0D2H1-2]
	\immini[thm]{Phần không bị gạch trong hình vẽ là miền nghiệm của bất phương trình nào sau đây?
		\choice
		{$y\geq x$}
		{$y\geq -x$}
		{$y\leq -x$}
		{\True $y\leq x$}}
	{\begin{tikzpicture}[scale=.7,thick,>=stealth,declare function={xmin=-2.5;xmax=3.5;ymin=-2;ymax=3;f(\x)=\x;}]
			\begin{scope}
				\clip (xmin,ymin) rectangle (xmax,ymax);
				\fill[pattern=north east lines] (xmin,ymin) rectangle (xmax,ymax);
				\fill[white] (xmin,{f(xmin)}) -- (xmax,{f(xmax)}) -- (xmax,ymin) -- (xmin,ymin) -- cycle;
				\draw[thick] (xmin,{f(xmin)}) -- (xmax,{f(xmax)});
			\end{scope}
			\draw[->] (xmin,0) -- (xmax,0)node[above]{$x$};
			\foreach \x in {-2,-1,1,2,3}
			\draw[shift={(\x,0)}] (0pt,2pt) -- (0pt,-2pt) node[below] {\footnotesize $\x$};
			\draw[->,color=black] (0,ymin) -- (0,ymax)node[left]{$y$};
			\foreach \y in {-1,1,2}
			\draw[shift={(0,\y)}] (2pt,0pt) -- (-2pt,0pt) node[left] {\footnotesize $\y$};
			\node[above left] at (0,0){$O$};
		\end{tikzpicture}}
	\loigiai{
		Đường thẳng bờ đi qua gốc toạ độ $O$ và điểm $(1;1)$ nên có phương trình $y=x$ (vẽ nét liền nên kể cả bờ).\\
		Lấy điểm $M(1;0)$ thuộc phần không bị gạch, ta có $0\leq 1$ nên phần không bị gạch là miền nghiệm của bất phương trình $y\leq x$.
	}
\end{ex}

%%%==============Câu_11 — y = 2x (2x - y = 0), tô nửa trên - trái, bờ NÉT LIỀN=============%%%
\begin{ex}%[VNPmath - tự biên soạn]%[GV. Vũ Ngọc Phát]%[0D2H1-2]
	\immini[thm]{Phần không bị gạch trong hình vẽ là miền nghiệm của bất phương trình nào sau đây?
		\choice
		{$y\geq 2x$}
		{$2y\geq x$}
		{$2y\leq x$}
		{\True $y\leq 2x$}}
	{\begin{tikzpicture}[scale=.7,thick,>=stealth,declare function={xmin=-2.5;xmax=3.5;ymin=-1.5;ymax=3.5;f(\x)=2*(\x);}]
			\begin{scope}
				\clip (xmin,ymin) rectangle (xmax,ymax);
				\fill[pattern=north east lines] (xmin,ymin) rectangle (xmax,ymax);
				\fill[white] (xmin,{f(xmin)}) -- (xmax,{f(xmax)}) -- (xmax,ymin) -- (xmin,ymin) -- cycle;
				\draw[thick] (xmin,{f(xmin)}) -- (xmax,{f(xmax)});
			\end{scope}
			\draw[->] (xmin,0) -- (xmax,0)node[above]{$x$};
			\foreach \x in {-2,-1,1,2,3}
			\draw[shift={(\x,0)}] (0pt,2pt) -- (0pt,-2pt) node[below] {\footnotesize $\x$};
			\draw[->,color=black] (0,ymin) -- (0,ymax)node[left]{$y$};
			\foreach \y in {-1,1,2,3}
			\draw[shift={(0,\y)}] (2pt,0pt) -- (-2pt,0pt) node[left] {\footnotesize $\y$};
			\node[above left] at (0,0){$O$};
		\end{tikzpicture}}
	\loigiai{
		Đường thẳng bờ đi qua gốc toạ độ $O$ và điểm $(1;2)$ nên có phương trình $y=2x$ (vẽ nét liền nên kể cả bờ).\\
		Lấy điểm $M(1;0)$ thuộc phần không bị gạch, ta có $0\leq 2$ nên phần không bị gạch là miền nghiệm của bất phương trình $y\leq 2x$.
	}
\end{ex}

%%%==============Câu_12 — y = x/2 + 1 (x - 2y + 2 = 0), tô nửa dưới - phải, bờ NÉT ĐỨT=============%%%
\begin{ex}%[VNPmath - tự biên soạn]%[GV. Vũ Ngọc Phát]%[0D2H1-2]
	\immini[thm]{Phần không bị gạch (không kể đường thẳng $d$) trong hình vẽ là miền nghiệm của bất phương trình nào sau đây?
		\choice
		{\True $x-2y+2<0$}
		{$x-2y+2>0$}
		{$2x-y+2<0$}
		{$2x-y+2>0$}}
	{\begin{tikzpicture}[scale=.7,thick,>=stealth,declare function={xmin=-3;xmax=3;ymin=-1.5;ymax=3.5;f(\x)=1+(\x)/2;}]
			\begin{scope}
				\clip (xmin,ymin) rectangle (xmax,ymax);
				\fill[pattern=north east lines] (xmin,ymin) rectangle (xmax,ymax);
				\fill[white] (xmin,{f(xmin)}) -- (xmax,{f(xmax)}) -- (xmax,ymax) -- (xmin,ymax) -- cycle;
				\draw[thick,dashed] (xmin,{f(xmin)}) -- (xmax,{f(xmax)});
			\end{scope}
			\draw[->] (xmin,0) -- (xmax,0)node[above]{$x$};
			\foreach \x in {-2,-1,1,2,3}
			\draw[shift={(\x,0)}] (0pt,2pt) -- (0pt,-2pt) node[below] {\footnotesize $\x$};
			\draw[->,color=black] (0,ymin) -- (0,ymax)node[left]{$y$};
			\foreach \y in {-1,1,2,3}
			\draw[shift={(0,\y)}] (2pt,0pt) -- (-2pt,0pt) node[left] {\footnotesize $\y$};
			\node[above left] at (0,0){$O$};
		\end{tikzpicture}}
	\loigiai{
		Đường thẳng bờ $d$ đi qua hai điểm $(0;1)$ và $(2;2)$ nên có phương trình $y=\dfrac{x}{2}+1$, hay $x-2y+2=0$.\\
		$d$ vẽ nét đứt nên không kể bờ.\\
		Lấy điểm $M(0;2)$ thuộc phần không bị gạch, ta có $0-2\cdot 2+2=-2<0$ nên phần không bị gạch là miền nghiệm của bất phương trình $x-2y+2<0$.
	}
\end{ex}

%%%==============Câu_13 — y = x + 2 (x - y + 2 = 0), tô nửa trên - trái, bờ NÉT ĐỨT=============%%%
\begin{ex}%[VNPmath - tự biên soạn]%[GV. Vũ Ngọc Phát]%[0D2H1-2]
	\immini[thm]{Phần không bị gạch (không kể đường thẳng $d$) trong hình vẽ là miền nghiệm của bất phương trình nào sau đây?
		\choice
		{$x-y+2<0$}
		{\True $x-y+2>0$}
		{$x+y+2<0$}
		{$x+y+2>0$}}
	{\begin{tikzpicture}[scale=.7,thick,>=stealth,declare function={xmin=-4;xmax=2;ymin=-1.5;ymax=3.5;f(\x)=(\x)+2;}]
			\begin{scope}
				\clip (xmin,ymin) rectangle (xmax,ymax);
				\fill[pattern=north east lines] (xmin,ymin) rectangle (xmax,ymax);
				\fill[white] (xmin,{f(xmin)}) -- (xmax,{f(xmax)}) -- (xmax,ymin) -- (xmin,ymin) -- cycle;
				\draw[thick,dashed] (xmin,{f(xmin)}) -- (xmax,{f(xmax)});
			\end{scope}
			\draw[->] (xmin,0) -- (xmax,0)node[above]{$x$};
			\foreach \x in {-3,-2,-1,1,2}
			\draw[shift={(\x,0)}] (0pt,2pt) -- (0pt,-2pt) node[below] {\footnotesize $\x$};
			\draw[->,color=black] (0,ymin) -- (0,ymax)node[left]{$y$};
			\foreach \y in {-1,1,2,3}
			\draw[shift={(0,\y)}] (2pt,0pt) -- (-2pt,0pt) node[left] {\footnotesize $\y$};
			\node[above left] at (0,0){$O$};
		\end{tikzpicture}}
	\loigiai{
		Đường thẳng bờ $d$ đi qua hai điểm $(-2;0)$ và $(0;2)$ nên có phương trình $y=x+2$, hay $x-y+2=0$.\\
		$d$ vẽ nét đứt nên không kể bờ.\\
		Lấy điểm $M(0;1)$ thuộc phần không bị gạch, ta có $0-1+2=1>0$ nên phần không bị gạch là miền nghiệm của bất phương trình $x-y+2>0$.
	}
\end{ex}

%%%%%%===================================================================%%%%%%
%%%%%%============================= DẠNG 4 ===============================%%%%%%
%%%%%% Bờ là đường thẳng ax+by = c với ab > 0  (hệ số góc âm)
%%%%%%===================================================================%%%%%%
\subsubsection{Bờ là đường thẳng $ax+by=c$ với $ab>0$}

%%%==============Câu_14 — x + y = 1, tô nửa dưới - trái, bờ NÉT ĐỨT=============%%%
\begin{ex}%[VNPmath - tự biên soạn]%[GV. Vũ Ngọc Phát]%[0D2H1-2]
	\immini[thm]{Phần không bị gạch (không kể đường thẳng $d$) trong hình vẽ là miền nghiệm của bất phương trình nào sau đây?
		\choice
		{$x+y<1$}
		{$x-y>1$}
		{\True $x+y>1$}
		{$x-y<1$}}
	{\begin{tikzpicture}[scale=.7,thick,>=stealth,declare function={xmin=-2.5;xmax=3.5;ymin=-2;ymax=3;f(\x)=1-(\x);}]
			\begin{scope}
				\clip (xmin,ymin) rectangle (xmax,ymax);
				\fill[pattern=north east lines] (xmin,ymin) rectangle (xmax,ymax);
				\fill[white] (xmin,{f(xmin)}) -- (xmax,{f(xmax)}) -- (xmax,ymax) -- (xmin,ymax) -- cycle;
				\draw[thick,dashed] (xmin,{f(xmin)}) -- (xmax,{f(xmax)});
			\end{scope}
			\draw[->] (xmin,0) -- (xmax,0)node[above]{$x$};
			\foreach \x in {-2,-1,1,2,3}
			\draw[shift={(\x,0)}] (0pt,2pt) -- (0pt,-2pt) node[below] {\footnotesize $\x$};
			\draw[->,color=black] (0,ymin) -- (0,ymax)node[left]{$y$};
			\foreach \y in {-1,1,2}
			\draw[shift={(0,\y)}] (2pt,0pt) -- (-2pt,0pt) node[left] {\footnotesize $\y$};
			\node[above left] at (0,0){$O$};
		\end{tikzpicture}}
	\loigiai{
		Đường thẳng bờ $d$ đi qua hai điểm $(1;0)$ và $(0;1)$ nên có phương trình $x+y=1$.\\
		$d$ vẽ nét đứt nên không kể bờ.\\
		Lấy điểm $M(2;0)$ thuộc phần không bị gạch, ta có $2+0=2>1$ nên phần không bị gạch là miền nghiệm của bất phương trình $x+y>1$.
	}
\end{ex}

%%%==============Câu_15 — x + y = 3, tô nửa dưới - trái, bờ NÉT LIỀN=============%%%
\begin{ex}%[VNPmath - tự biên soạn]%[GV. Vũ Ngọc Phát]%[0D2H1-2]
	\immini[thm]{Phần không bị gạch trong hình vẽ là miền nghiệm của bất phương trình nào sau đây?
		\choice
		{$x+y<3$}
		{\True $x+y\geq 3$}
		{$x-y\geq 3$}
		{$x+y\leq 3$}}
	{\begin{tikzpicture}[scale=.7,thick,>=stealth,declare function={xmin=-1.5;xmax=4.5;ymin=-1;ymax=4;f(\x)=3-(\x);}]
			\begin{scope}
				\clip (xmin,ymin) rectangle (xmax,ymax);
				\fill[pattern=north east lines] (xmin,ymin) rectangle (xmax,ymax);
				\fill[white] (xmin,{f(xmin)}) -- (xmax,{f(xmax)}) -- (xmax,ymax) -- (xmin,ymax) -- cycle;
				\draw[thick] (xmin,{f(xmin)}) -- (xmax,{f(xmax)});
			\end{scope}
			\draw[->] (xmin,0) -- (xmax,0)node[above]{$x$};
			\foreach \x in {-1,1,2,3,4}
			\draw[shift={(\x,0)}] (0pt,2pt) -- (0pt,-2pt) node[below] {\footnotesize $\x$};
			\draw[->,color=black] (0,ymin) -- (0,ymax)node[left]{$y$};
			\foreach \y in {-1,1,2,3}
			\draw[shift={(0,\y)}] (2pt,0pt) -- (-2pt,0pt) node[left] {\footnotesize $\y$};
			\node[above left] at (0,0){$O$};
		\end{tikzpicture}}
	\loigiai{
		Đường thẳng bờ đi qua hai điểm $(3;0)$ và $(0;3)$ nên có phương trình $x+y=3$ (vẽ nét liền nên kể cả bờ).\\
		Lấy điểm $M(4;0)$ thuộc phần không bị gạch, ta có $4+0=4\geq 3$ nên phần không bị gạch là miền nghiệm của bất phương trình $x+y\geq 3$.
	}
\end{ex}

%%%==============Câu_16 — x + 2y = 2, tô nửa trên - phải, bờ NÉT ĐỨT=============%%%
\begin{ex}%[VNPmath - tự biên soạn]%[GV. Vũ Ngọc Phát]%[0D2H1-2]
	\immini[thm]{Phần không bị gạch (không kể đường thẳng $d$) trong hình vẽ là miền nghiệm của bất phương trình nào sau đây?
		\choice
		{$x+2y>2$}
		{$2x+y>2$}
		{\True $x+2y<2$}
		{$2x+y<2$}}
	{\begin{tikzpicture}[scale=.7,thick,>=stealth,declare function={xmin=-2;xmax=4;ymin=-2;ymax=3;f(\x)=1-(\x)/2;}]
			\begin{scope}
				\clip (xmin,ymin) rectangle (xmax,ymax);
				\fill[pattern=north east lines] (xmin,ymin) rectangle (xmax,ymax);
				\fill[white] (xmin,{f(xmin)}) -- (xmax,{f(xmax)}) -- (xmax,ymin) -- (xmin,ymin) -- cycle;
				\draw[thick,dashed] (xmin,{f(xmin)}) -- (xmax,{f(xmax)});
			\end{scope}
			\draw[->] (xmin,0) -- (xmax,0)node[above]{$x$};
			\foreach \x in {-1,1,2,3,4}
			\draw[shift={(\x,0)}] (0pt,2pt) -- (0pt,-2pt) node[below] {\footnotesize $\x$};
			\draw[->,color=black] (0,ymin) -- (0,ymax)node[left]{$y$};
			\foreach \y in {-1,1,2}
			\draw[shift={(0,\y)}] (2pt,0pt) -- (-2pt,0pt) node[left] {\footnotesize $\y$};
			\node[above left] at (0,0){$O$};
		\end{tikzpicture}}
	\loigiai{
		Đường thẳng bờ $d$ đi qua hai điểm $(0;1)$ và $(2;0)$ nên có phương trình $y=1-\dfrac{x}{2}$, hay $x+2y=2$.\\
		$d$ vẽ nét đứt nên không kể bờ.\\
		Lấy điểm $O(0;0)$ thuộc phần không bị gạch, ta có $0+2\cdot 0=0<2$ nên phần không bị gạch là miền nghiệm của bất phương trình $x+2y<2$.
	}
\end{ex}


\Closesolutionfile{ans}
